%!TEX root = ./../../springer-mono-format/book.tex
%!TEX root = ./../../thesis-trier.tex

Here, we investigate the following subclasses of automata with respect to the notions
of complete distinguishability, complete reachability,
and sync-maximality: solvable automata (\Cref{sync_max:sec:blocks}) 
commutative automata (\Cref{sync_max:sec:comm}), automata where every letter has defect at most one (\Cref{sync_max:sec:defectone}), SI-automata (Section~\ref{sync_max:sec:SI}),
circular automata (Section~\ref{sync_max:sec:circular}), where \cref{sync_max:sec:circular} 
also introduces circular automata, and binary automata (\Cref{sync_max:sec:binary}).

\section{Automata Whose Strongly Connected Components are Blocks}
\label{sync_max:sec:blocks}

%!TEX root = ./../../springer-mono-format/book.tex
%!TEX root = ./../../thesis-trier.tex

In an automaton whose strongly connected components are blocks, a synchronizing state
is a trap state by \cref{prelim:class:block:cor:sync_trap_state}.
Thus, by \cref{prelim:subclasses:sync:cor:alg-unique-trap} for these automata
we can decide synchronizability in linear time and constant space.
Furthermore, from \cref{prelim:lem:compl_reach_implies_SCC_sync,sync_max:thm:char_sync_max_aut} 
we obtain that such an automaton is completely reachable if and only if it has a single state,
and sync-maximal if and only if it has at most two states and is synchronizing with a trap state.
Recall that this class includes partially ordered, commutative,
and solvable automata (for partially 
ordered automata the time bound on deciding synchronizability was given in~\cite[Cor.~6]{DBLP:conf/dlt/Hoffmann21a}).

\begin{theorem} \label{syncmax:thm:solvable-cd}
	Let $\semiautomat{A}{Q}{\Sigma}{\delta}$ be an automaton in which every strongly connected component is a block.
	If~$\mathcal A$ is completely distinguishable,
	then $|Q| \le 2$.		
\end{theorem}
\begin{proof}
	By \cref{constr_sync:cor:solvable_sync_iff_reach_trap}, $\mathcal A$ is synchronizing
	and by \cref{prelim:class:block:cor:sync_trap_state} the synchronizing
	state must be a trap state~$t \in Q$. %, which is reachable from every other state by \cref{prelim:cor:sync_trap_state}.
	
	Since $t$ is not the only state of~$\mathcal A$,
	we can choose some maximal $p \in Q$ w.r.t.~$\lereach$
	that is not equal to~$t$.
	This is equivalent to the condition
	that, for each $x \in \Sigma$, we have $\delta(p, x) = t$
	or $\delta(p, x) \eqreach p$. Note that in the latter case,
	$\delta(p,x)$ is also maximal and distinct from~$t$.
	
	Now, assume there exists another $q \in Q \setminus \{p,t\}$
	with this property.
	Then the sets $\{p,q\}$ and $\{p,q,t\}$ are indistinguishable.
	For assume $u = u_1 \cdots u_m$ with $u_1, \ldots, u_m \in \Sigma$
	distinguishes both sets.
	First, there must exist a minimal $i \in \{1,\ldots,m\}$
	such that, without loss of generality,
	$\delta(p, v_i) = t$ for $v_i = u_1 \cdots u_i$. 
	For assume this is not the case.
	Then, inductively, $\delta(p, v_i) \eqreach p$
	and $\delta(q, v_i) \eqreach q$,
	which also implies that none of these states equals~$t$.
	Since each strongly connected component is a block,
	we also have $\delta(p, v_i) \ne \delta(q, v_i)$.
	Thus $u$ cannot distinguish both sets.
	However, $\delta(p, v_i) = t$
	implies  $\delta(\{p,q\}, v_i) = \delta(\{p,q,t\}, v_i)$,
	and so~$u$ cannot distinguish both sets in this case as well.
	Thus, there exists no state $q$ as above, which 
	implies that~$\{p\}$ is a strongly connected
	component and that every path from every state distinct from~$t$ to~$t$
	has to pass through the state~$p$.
	
	Now, let $q \in Q$ be a state distinct from $p$ and~$t$
	that is minimal w.r.t.~$\lereach$.
	A word distinguishing both sets has to map $\{q,t\}$
	to a singleton, since it is a subset of the other set,
	which must be $\{t\}$, since~$t$ is a trap state.
	Now, since a prefix of a word mapping~$q$ to~$t$ has to map~$q$ to~$p$,
	this prefix maps both sets to the same set
	and thus cannot distinguish both sets.
	This yields that both sets are indistinguishable.
\end{proof}

Recall that, by \cref{syncmax:rem:two-state-compl-dist}, automata with at most two states are 
always completely distinguishable.

\section{Commutative Automata}
\label{sync_max:sec:comm}

%!TEX root = ./../../springer-mono-format/book.tex
%!TEX root = ./../../thesis-trier.tex

As follows \cref{sync_max:sec:blocks},
a commutative automaton that is completely reachable
has a single state, and if it completely distinguishable
it has at most two states.\footnote{By \cref{prelim:classes:comm:prop:permSCC-regular},
	the transformation monoid of a strongly connected commutative $n$-state	automaton has size~$n$.	
	But as noted in \cref{sync_max:sec:compl_reach} the transformation monoid of a completely reachable automaton
	has size at least $2^{|Q|} - |Q|$. This also immediately implies $|Q| \le 2$ for completely reachable commutative automata.}

Next, for a fixed alphabet, we give polynomial-time algorithms to compute (i) a partially ordered automaton accepting the intersection of partially ordered commutative automata (\cref{compl:prop:intersection_waa_com}) and (ii) an automaton accepting the set of synchronizing words of a commutative automaton (\cref{prop:aut_for_sync_words}). The latter result implies that there is no exponential blow-up in the size of automata accepting this set. 
As consequences, stated in \cref{compl:cor:intersection_waa_com,cor:aut_for_sync_words},
the problems \DFAINT and \DFASWS are solvable in polynomial time for a fixed alphabet.
Lastly, in \cref{syncmax:prop:unary-sync-words} we give an expression for the set of synchronizing words of a unary automaton.

\Cref{compl:prop:intersection_waa_com,prop:aut_for_sync_words} 
were first presented in~\cite[Prop.~41 \& Lem.~43]{DBLP:journals/tcs/Hoffmann21} by the author.

\begin{proposition}
	\label{compl:prop:intersection_waa_com}
	Let $\mathcal A_i = (Q_i,\Sigma, \delta_i, q_i, F_i)$,
	$i \in \{1,\ldots,m\}$,
	be partially ordered commutative automata
	with at most $n$ states.
	Then, $\bigcap_{i=1}^m L(\mathcal A_i)$
	is recognizable by a partially ordered
	commutative automaton 
	of size $n^{|\Sigma|}$ % at most
	computable in time $O(m|\Sigma| n^{|\Sigma|+1})$.
\end{proposition} 
\begin{proof} 
	Fix an ordering $\Sigma = \{a_1, \ldots, a_k\}$ with $k \ge 1$.
	By~\cref{prelim:classes:cor:com-po} we obtain
	\begin{equation}\label{compl:eqn:com_pto_int}
		u \in \bigcap_{i=1}^m L(\mathcal A_i)
		\Leftrightarrow a_1^{\min\{n-1, |u|_{a_1}\}} \cdots a_k^{\min\{n-1, |u|_{a_k}\}} \in \bigcap_{i=1}^m L(\mathcal A_i)
	\end{equation}
	as $n \ge \max\{ |Q_i| : i \in \{1,\ldots,m\} \} - 1$. 
	Next, set $\automat{A}{Q}{\Sigma}{\delta}{q_0}{F}$
	where $Q = [n]^k$, $q_0 = (0,\ldots,0)$,
	\[
	\applytrans{\delta}{(s_1, \ldots, s_k)}{a_j} = (s_1, \ldots, s_{j-1}, \min\{n-1, s_j+1\}, s_{j+1}, \ldots, s_k),
	\]
	and 
	\[
	F = \left\{ (s_1, \ldots, s_k) \in Q \mid a_1^{s_1} \cdots a_k^{s_k} \in \bigcap_{i=1}^m L(\mathcal A_i)
	\right\}.
	\]
	Then we have
	\begin{align*}
		u \in \bigcap_{i=1}^m L(\mathcal A_i)
		& \Leftrightarrow  a_1^{\min\{n-1, |u|_{a_1}\}} \cdots a_k^{\min\{n-1, |u|_{a_k}\}} \in \bigcap_{i=1}^m L(\mathcal A_i) 
		& \tag{Eqn.~\eqref{compl:eqn:com_pto_int}}\\
		& \Leftrightarrow (\min\{n-1, |u|_{a_1}\}, \ldots, \min\{n-1, |u|_{a_k}\}) \in F 
		& \tag{Def.~$F$} \\ 
		& \Leftrightarrow \applytrans{\delta}{q_0}{a_1^{|u|_{a_1}} \cdots a_k^{|u|_{a_k}}} \in F 
		& \tag{Def.~$\delta$} \\
		& \Leftrightarrow a_1^{|u|_{a_1}}	\cdots a_k^{|u|_{a_k}} \in L(\mathcal A) 
		& \tag{Def. $L(\mathcal A)$}\\
		& \Leftrightarrow u \in L(\mathcal A). \tag{$\mathcal A$ commutative}					
	\end{align*}
	So $\bigcap_{i=1}^m L(\mathcal A_i) = L(\mathcal A)$.
		
	The automaton, without~$F$, can clearly
	be computed in time $O(n^k)$. To
	compute~$F$, 
	iterate through all $n^k$ words $a_1^{s_1} \cdots a_k^{s_k}$
	with $(s_1, \ldots, s_k) \in [n]^k$
	and test for each word whether
	$a_1^{s_1} \cdots a_k^{s_k} \in L(\mathcal A_i)$
	for all automata $\mathcal A_i$, $i \in \{1,\ldots,m\}$.
	For $i \in \{1,\ldots,m\}$, testing a single word  $a_1^{s_1} \cdots a_k^{s_k}$ of length in~$O(kn)$
	whether it is accepted by $\mathcal A_i$
	can be done in time~$O(kn)$. So testing 
	whether it is accepted by all automata $\mathcal A_i$, $i\in \{1,\ldots,m\}$,
	can be done in time $O(mkn)$. 
	In total, the algorithm runs in time $O(mkn^{k+1})$. 
\end{proof}

\begin{remark}
	The bound from \cref{compl:prop:intersection_waa_com} is tight.
	For example, set $\Sigma = \{a_1, \ldots, a_m\}$
	and for $i \in \{1,\ldots,m\}$ consider $\mathcal A_i = ([n], \Sigma, \delta_i, 0, \{n-1\})$
	with $\delta(j, a_i) = \min\{j+1,n-1\}$ and $\delta(j, a_k) = j$
	for $k \ne i$ and $j \in [n]$. We have $L(\mathcal A_i) = \{ u \in \Sigma^* : |u|_{a_i} \ge n\}$
	and $L = \bigcap_{i=1}^m L(\mathcal A_i) =  \{ u \in \Sigma^* : \forall i \in \{1,\ldots,k\} \colon |u|_{a_i} \ge n \}$.
	Let $0 \le i_1, \ldots, i_k, j_1, \ldots, j_k < n$
	where $i_r < j_r$ for some $r \in \{1,\ldots,k\}$.
	Set $u = a_1^{i_1} \cdots a_k^{i_k}$,
	$v = a_1^{j_1} \cdots a_k^{j_k}$,	
	and $w = a_1^{n - j_1} \cdots a_k^{n - j_k}$.
	Then
	\[
	uw \notin L \text{ and }
	vw \in L.
	\]
	Thus $u \not\equiv_L v$.
	This yields that the words $a_1^{i_1} \cdots a_k^{i_k}$
	for $0 \le i_1, \ldots, i_k < n$
	induce distinct Myhill-Nerode equivalence classes.
	Then every automaton accepting~$L$ has at least $n^k$ states
	and thus the bound from \cref{compl:prop:intersection_waa_com}
	is reached.	
	This also follows from results in~\cite{DBLP:journals/iandc/Hoffmann24}
	where the state complexity on commutative languages is investigated.
	The above languages were used in the proof of \cite[Thm.~1]{DBLP:journals/ijfcs/BrzozowskiL12},
	so that this proof established the above in the special case $m = 2$.
\end{remark}

\begin{corollary}	\label{compl:cor:intersection_waa_com}
	\DFAINT is solvable in time $O(m|\Sigma|n^{|\Sigma|+1})$
	for $m$ partially ordered input automata each having at most $n$ states.	
\end{corollary}
\begin{proof}
	The input automata accept a common word if and only if 
	the final state set in the automaton constructed
	in the proof of~\cref{compl:prop:intersection_waa_com}
	is non-empty.
	So, we can use the same algorithm
	as in the proof of~\cref{compl:prop:intersection_waa_com}
	and simply test whether $F \ne \emptyset$
	at the end.
\end{proof}

\begin{proposition}
	\label{prop:aut_for_sync_words}
	Let $\semiautomat{A}{Q}{\Sigma}{\delta}$
	be a commutative automaton with $n$ states.
	Then, the set of synchronizing words
	is recognizable by a partially ordered commutative automaton
	of size~$n^{|\Sigma|}$
 	computable in time $O(|\Sigma|n^{|\Sigma|+2})$.
\end{proposition}
\begin{proof}		
	By \cref{prelim:classes:comm:cor:collections-of-facts:2}, 
	for $\Syn(\mathcal A) \ne \emptyset$,
	$\mathcal A$ must have a unique trap state.
	So, first check whether $\mathcal A$
	has a unique trap state $t \in Q$, and compute it if present, 
	which can be done in time $O(n)$ by \cref{prelim:subclasses:sync:alg-unique-trap}.
	In case $\mathcal A$
	has no trap state, then $\Syn(\mathcal A) = \emptyset$
	so that a one-state automaton with empty final state set
	accepts $\Syn(\mathcal A)$.
	Otherwise,  as, by \cref{constr_sync:prop:pto_comm_rho_automata},~$\mathcal A$ is a $\rho$-automaton
	with a trap state,
	we find $\Syn(\mathcal A) = \Syn(\mathcal A / {\eqreach})$  by \cref{constr_input:prop:rho_trap_sync_factor}.
	As a quotient automaton of a commutative automaton,
	$\mathcal A / {\eqreach} = (Q / {\eqreach}, \Sigma, \delta_{\eqreach})$
	is commutative, and has the trap state $[t]_{\eqreach} = \{t\}$.
	By \cref{constr_input:prop:factor_rho_aut_pto}, 
	$\mathcal A / {\eqreach}$
	is a partially ordered automaton
	computable in linear time.	
	As, by \crefitem{prelim:thm:sync_trap_state}{prelim:subclasses:thm:trap:item:syncunique},~$[t]_{\eqreach}$ must be the unique synchronizing state, we have
	\[
	\Syn(\mathcal A / {\eqreach}) = \bigcap_{[q]_{\eqreach} \in Q / {\eqreach}} L(( \mathcal A / {\eqreach} )_{[q]_{\eqreach}, \{[t]_{\eqreach}\}}).
	\]
	By \cref{compl:prop:intersection_waa_com} 
	we can compute a partially ordered commutative
	automaton of size $n^{|\Sigma|}$ accepting this intersection involving $n$ automata,
	each having~$n$ states,	in time $O(|\Sigma|n^{|\Sigma|+2})$.
	
	The above intersection for $\Syn(\mathcal A / {\eqreach})$
	equals the same intersection but only over
	the minimal states $Q / {\eqreach}$, which might give a slightly better running time for the algorithm.
\end{proof}

We note a direct consequence of the previous result.

\begin{corollary}	\label{cor:aut_for_sync_words}
	\DFASWS\ is solvable in time $O(|\Sigma|n^{|\Sigma|+2})$
	for $n$-state commutative input automata $\semiautomat{A}{Q}{\Sigma}{\delta}$.	
\end{corollary}
\begin{proof} 
	By \cref{prop:aut_for_sync_words}, 
	we can compute an automaton of size $n^{|\Sigma|}$
	for $\Syn(\mathcal A)$ in time $O(|\Sigma|n^{|\Sigma|+2})$. Then a shortest word accepted
	by this automaton can be computed in time $O(|\Sigma|n^{|\Sigma|})$
	as a shortest path from the start state to some final state~\cite{cormen01introduction}.
\end{proof}

Martyugin~\cite[Prop.~5.2]{Mar2009b} provided a direct algorithm to decide \DFASWS  
running in time $O(|\Sigma|n^{|\Sigma|}\log n)$ for commutative input automata 
(he actually omits the factor~$|\Sigma|$ as he assumed a fixed alphabet; 
however, the proof shows that this is the running time 
when the alphabet size is taken into account).

Lastly, we note the following improvement for unary languages.

\begin{proposition} \label{syncmax:prop:unary-sync-words}
	Let $\semiautomat{A}{Q}{\{a\}}{\delta}$
	be a synchronizing unary automaton.
	Then $\Syn(\mathcal A) = a^{l} a^*$
	where $l$ is the length of a longest path traversing distinct states in~$\mathcal A$.
	In particular, $\Syn(\mathcal A) = a^{l} a^*$
	is acceptable by an automaton with at most~$l + 1\le |Q|$ states.
\end{proposition}
\begin{proof}
	A synchronizing unary automaton consists of tails all ending at a common state
	at which the single letter induces a self-loop.
	Now, it is easy to see that the set of synchronizing words
	has the form as claimed. Clearly, $l < |Q|$
	and $a^{l} a^*$ is acceptable by an automaton with $l + 1$ states.
\end{proof}

The smallest automaton whose set of synchronizing words is $a^la^*$
equals the minimal automaton accepting this language over $\{a\}$.
Thus the reset complexity and the state complexity of $a^la^*$ both equal $l+1$.
This was also stated in~\cite[Prop.~1]{DBLP:journals/corr/Maslennikova14}.

\section{All Letters Have Defect at Most One}
\label{sync_max:sec:defectone}

%!TEX root = ./../../springer-mono-format/book.tex
%!TEX root = ./../../thesis-trier.tex

Here, \cref{thm:defectatmostone}
shows that for automata where every letter has defect at most one
and whose transformation monoid contains
a transitive permutation group, synchronizability implies complete distinguishability.
\Cref{cor:defectatmostone-weaker} lists some equivalences that follow
from \cref{thm:defectatmostone}.
Lastly, in \cref{cor:single-perm,cor:char-bin-compl-reach,cor:bin-compl-reach-sync-maximal},
we conclude that in automata with a single non-permutational letter	
and for binary automata, complete reachability
and sync-maximality are equivalent properties,
and that every binary completely reachable is sync-maximal.
The latter statement answers the open question
from~\cite{HoffmannCircularExtended} whether there exist 
binary completely reachable automata that are not sync-maximal.

\begin{theorem}\label{thm:defectatmostone}
	Let $\semiautomat{A}{Q}{\Sigma}{\delta}$
	be an automaton such that the group of units of~$\TM{\mathcal A}$
	is transitive and every letter has defect at most one.
	If $\mathcal A$ is synchronizing, then it is completely distinguishable.	
\end{theorem}
\begin{proof}
	Since all automata with at most two states are completely distinguishable by \cref{syncmax:rem:two-state-compl-dist},
	we can assume $|Q| \ge 3$.
	Let $u \in \Sigma^*$ be a shortest word such that
	$\applytrans{\delta}{Q}{u} = \{q\}$ for some $q \in Q$.
	Then we can write $u = xv$  with  $x \in \Sigma$ and $v \in \Sigma^*$.
	By the minimal choice of~$u$, we have $|\applytrans{\delta}{Q}{v}| \ge 2$
	and $\delta(Q, x) \ne Q$.
	By assumption,~$x$ has defect one,
	so that $\applytrans{\delta}{Q}{x} = Q \setminus \{p\}$
	for some state $p \in Q$.
	Hence, as $\delta(Q \setminus \{p\}, v) = \{q\}$, we get
	\begin{align*}
		\applytrans{\delta}{Q}{v}  = \{\applytrans{\delta}{p}{v} \} \cup \applytrans{\delta}{Q \setminus \{p\}}{v} 
		= \{ \applytrans{\delta}{p}{v}, q \}
	\end{align*}
	where $\applytrans{\delta}{p}{v} \ne q$.    
	Now, let $r \in Q$.	
	As $\TM{\mathcal A}$ contains a transitive permutation group, 
	there exists $w \in \Sigma^*$ inducing a permutation on $Q$
	such that $\delta(r, w) = p$.
	Then 
	$\applytrans{\delta}{Q \setminus \{r\}}{wv} = \applytrans{\delta}{Q\setminus\{p\}}{v} = \{q\}$
	and $\applytrans{\delta}{r}{wv} = \applytrans{\delta}{p}{v}$.
	As $\applytrans{\delta}{r}{wv} \notin \applytrans{\delta}{Q \setminus\{r\}}{wv}$
	and $|\applytrans{\delta}{Q\setminus \{r\}}{wv}| = 1$,
	by~\crefitem{sync_max:thm:char_distinguishable}{dist:item:single_point},
	it follows that $\mathcal A$ is completely distinguishable.
\end{proof}

\begin{remark}
	The automaton from \cref{sync_max:ex:2-sets-not-all-dist},
	and those from \cref{sync_max:fig:sync-two-sets-dist}
	given later, only have simple idempotent letters (and so of defect one),
	are completely reachable, but they are not completely distinguishable.
	Thus, strong connectivity and synchronizability
	do not imply complete distinguishability for these automata.
\end{remark}


\begin{corollary} \label{cor:defectatmostone-weaker}
	Let $\semiautomat{A}{Q}{\Sigma}{\delta}$
	be an automaton such that the group of units of~$\TM{\mathcal A}$
	is transitive (equivalently $\restrictAutGamS{\mathcal A}{\Gamma}{Q}$
	for $\Gamma = \{ x \in \Sigma \mid \text{$x$ is a permutational letter } \}$ is strongly connected) and every letter has defect at most one with $|Q| \ge 5$.
	Then the following are equivalent.
	\begin{enumerate}
	\item \label{defone:item:sync} $\mathcal A$ is synchronizing.
	
	\item \label{defone:item:cd} $\mathcal A$ is completely distinguishable.
	
	\item \label{defone:item:k-sets} 	
	Every $k$-sets of states is distinguishable from every distinct $l$-set of states for some $2 \le k, l \le |Q|$ with $|Q| > k$ or $|Q| > l$.
		
	\item \label{defone:item:dist-card} For two distinct non-empty and non-singleton subsets $A, B \subseteq Q$
	we have $|\delta(A, u)| \ne |\delta(B, u)|$ for some $u \in \Sigma^*$.		
	
	\item \label{defone:item:disjoint-kl-sets} 
	Every $k$-set of states is distinguishable from every disjoint $l$-set
	of states for some $k, l \ge 2$ with $kl < |Q|$.
	\end{enumerate}
\end{corollary}
\begin{proof}
$\ref{defone:item:sync} \Rightarrow \ref{defone:item:cd}$:
	This follows from \cref{thm:defectatmostone}.\footnote{For SI-automata, this also follows from the graph-theoretical methods from \cref{sync_max:chap:graphs}.
For consider the Rystsov graph $I(\TM{\mathcal A})$.
By the assumption the group of units of $\TM{\mathcal A}$
is transitive.
As the group of units is part of the automorphism group 
of $I(\TM{\mathcal A})$ (Proposition~\ref{sync_max:prop:G_automorhism_on_graphs}) 
the graph $I(\TM{\mathcal A})$
is vertex-transitive.
As the completely reachable automaton $\mathcal A$ 
is synchronizing (Lemma~\ref{prelim:lem:compl_reach_implies_SCC_sync}),
by \cref{sync_max:thm:rystsov_SI} and \crefitem{sync_max:prop:Rystsov_graph_vs_SI_graph}{sync_max:prop:item:IAsubgraphIS} 
the graph $I(\TM{\mathcal A})$ is strongly connected.
By \cref{prelim:prop:vertex-transitive-SCC-32-connected} 
it is strongly $\frac{3}{2}$-connected.
So Theorem~\ref{sync_max:thm:IS_32SCC} yields that $\mathcal A$ is 
completely distinguishable.}

$\ref{defone:item:cd} \Rightarrow \ref{defone:item:k-sets}$: This is obvious.

$\ref{defone:item:k-sets} \Rightarrow \ref{defone:item:sync}$:
For $k \ge 3$ or $l \ge 3$, by \cref{sync_max:cor:k-ge-3-sync} the automaton $\mathcal A$ is synchronizing.
For $k = l = 2$, since by assumption $\mathcal A$ is strongly connected it follows from \cref{sync_max:prop:2-sets-dist-SCC-implies-sync} that $\mathcal A$ is synchronizing. 

\ref{defone:item:cd} $\Leftrightarrow$ \ref{defone:item:dist-card}: If~$\mathcal A$
is completely distinguishable, then precisely one subset from $A$ or $B$
is mapped to a singleton while the other is mapped to a set of size at least two.
Thus, \cref{defone:item:dist-card} holds. Conversely,
if \cref{defone:item:dist-card}, then distinct pairs of $2$-sets
are distinguishable and since the equivalence of \cref{defone:item:k-sets} 
and \cref{defone:item:cd} has already been established,
we obtain that the automaton is completely distinguishable.

\ref{defone:item:disjoint-kl-sets} $\Leftrightarrow$ \ref{defone:item:sync}: 
First, suppose Item~\ref{defone:item:disjoint-kl-sets} holds.
Then, \cref{sync_max:prop:disjointklsets} applied to $\TM{\mathcal A$} yields that~$\mathcal A$ is synchronizing.
Conversely, if~$\mathcal A$ is synchronizing, then since the other equivalences
have already been established it is completely distinguishable,
in particular the pairs of sets as stated are distinguishable.
\end{proof}

In general, distinguishability of pairs of distinct $2$-sets does not imply
synchronizability even for SI-automata, see \cref{sync_max:ex:2-sets}. For strongly connected 
automata, distinguishability of pairs of distinct $2$-sets does imply
synchronizability by \cref{sync_max:prop:2-sets-dist-SCC-implies-sync},
but it does not imply complete distinguishability even for strongly connected SI-automata,
see the examples from~\cref{sync_max:sec:furtherexamples}.
\Cref{sync_max:ex:2-sets-but-not-all-new,sync_max:ex:2-sets-not-all-dist} showed that, in general, the
ability to map two distinct non-singleton subsets to subsets of different cardinality
is not sufficient for complete distinguishability.

\begin{corollary} \label{cor:single-perm}
	Let $\semiautomat{A}{Q}{\Sigma}{\delta}$
	be an automaton 
	with at least three states and a single non-permutational letter.
	Then $\mathcal A$ is completely reachable
	if and only if it is sync-maximal.		
\end{corollary}
\begin{proof}
	First, suppose~$\mathcal A$ is completely reachable.
	By \cref{sync_max:cor:single-non-permutation}
	the group of units of $\TM{\mathcal A}$	is transitive
	and the single non-permutational letter has defect one.
	So, by \cref{thm:defectatmostone}, since $\mathcal A$ is synchronizing
	we find that it is completely distinguishable.
	From \crefitem{sync_maximal:thm:structure}{sync_maximal:thm:structure:1}
	we get that $\mathcal A$ is sync-maximal.
		
	Conversely, suppose~$\mathcal A$ is sync-maximal.
	If $\mathcal A$ is not completely reachable,
	then, by \cref{sync_maximal:thm:structure}, it is not strongly connected,
	and so, by \crefitem{sync_max:cor:structure_letters}{sync_max:cor:structure_letters:1},
	has two distinct non-permutational letters.
	Hence, as $\mathcal A$ has only one non-permutational letter, it must be completely reachable.		
\end{proof}

\begin{corollary} \label{cor:char-bin-compl-reach}
	Let $\semiautomat{A}{Q}{\Sigma}{\delta}$
	be a binary automaton 
	with at least three states.
	Then $\mathcal A$ is completely reachable
	if and only if it is sync-maximal.		
\end{corollary}
\begin{proof}
	First, suppose $\mathcal A$ is completely reachable. By \cref{sync_max:cor3:single-non-permutation}, $\mathcal A$
	has precisely one non-permutational letter.
	So, by \cref{cor:single-perm}{cor:single-perm:1}, it is sync-maximal.
	Conversely, by \cref{sync_maximal:thm:structure} and \cref{sync_max:cor:structure_letters}, 
	a binary sync-maximal
	automaton with at least three states
	is completely reachable.	
\end{proof}

\begin{corollary} \label{cor:bin-compl-reach-sync-maximal}
	Every binary completely reachable automaton is sync-maximal.
\end{corollary}
\begin{proof}
	This follows from \cref{cor:char-bin-compl-reach} 
	since by \cref{sync_max:cor:char_sync_max_aut} 
	synchronizing automata with at most two states are sync-maximal.
\end{proof}

\cref{cor:single-perm,cor:char-bin-compl-reach}
do not hold for two-state automata,
as by \crefitem{sync_max:cor:char_sync_max_aut}{cha_sync_max:item:Q2} 
every two-state synchronizing automaton is sync-maximal,
but there exist two-state synchronizing automata that are not strongly
connected, and so not completely reachable.

As complete reachability is decidable in polynomial time~\cite{DBLP:conf/icalp/FerensS23}, it follows from \cref{cor:char-bin-compl-reach} that, for binary automata, sync-maximality is also decidable in polynomial time. We note that Casas and Volkov~\cite{DBLP:conf/latin/CasasV22} gave a characterization of binary completely reachable automata (see also the discussion in~\cref{sync_max:sec:subclasses:discussion}), as well as a faster polynomial-time algorithm than the general one from~\cite{DBLP:conf/icalp/FerensS23} for deciding complete reachability in the binary case.

\section{Automata with Simple Idempotents}
\label{sync_max:sec:SI}

%!TEX root = ./../../thesis-trier.tex

First, in \cref{sec:SI:dist}, we prove
that in a synchronizing SI-automaton every pair of distinct $2$-sets is distinguishable.
Lastly, in \cref{sec:SI:SCC} we investigate strongly connected SI-automata 
and derive conditions connecting synchronizability, complete reachability,
and complete distinguishability. We also derive these connections for binary SI-automata in this section.

\subsection{Distinguishability of $2$-Sets}
\label{sec:SI:dist} 

%!TEX root = ./../../thesis-trier.tex

Here, in \cref{syncmax:cor:SI-2-sets}, we prove that in a synchronizing SI-automaton
all pairs of distinct $2$-sets are distinguishable.
This follows from a series of preliminary results 
regarding mapping properties of equal-sized sets.

\begin{lemma} \label{app:lem:SI}
	Let~$t \colon X \to X$ be an one-point singular
	and $A,B \subseteq X$ be finite
	with $|A| = |B|$ and $t(A) = t(B)$.
	Then $\move{t} \cap A = \move{t} \cap B$ implies $A = B$.
\end{lemma}
\begin{proof} 
	Let $x \in X$ be such that $\dupl{t} = \{x\}$.
	Consider $z \in A$ with $z \ne x$.
	If $t(z) \ne z$, then $z \in B$ by assumption.
	Next, suppose $t(z) = z$.
	Since $t(A) = t(B)$,
	we have $t(z) = t(y)$ for some $y \in B$.
	As $t$ is a one-point singular
	and $t(z) \ne x$,
	we find $t^{-1}(t(z)) = \{z\}$. Therefore $y = z$
	and so $z \in B$.
	This gives $A \setminus \{x\} \subseteq B$.
	Similarly, we obtain the other inclusion
	and thus $A \setminus \{x\} = B \setminus \{x\}$.
	Since $|A| = |B|$, 
	either $x \in A \cap B$ or $x \notin A \cup B$.
	This gives $A = B$.	
\end{proof}

\begin{corollary} \label{app:cor:SI}
		Let $t \colon X\to X$ be a simple idempotent
		and $A, B \subseteq X$ be finite.
		Suppose $|A| = |B|$, $\applyfunc{A}{t} = \applyfunc{B}{t}$, 
		$\applyfunc{A}{t} \ne A$, and $\applyfunc{B}{t} \ne B$.
		Then $A = B$.	
\end{corollary}
\begin{proof}
Let $\move{t} = \{a\}$.
Since $t(A) \ne A$ and $t(B) \ne B$, we have $a \in A \cap B$.	
Thus, the claim follows from \cref{app:lem:SI}.
\end{proof}

By \cref{app:lem:SI}, if $A \ne B$ with $|A| = |B|$ and $\applyfunc{A}{t} = \applyfunc{B}{t}$,
then, without loss of generality, $\applyfunc{B}{t} \ne B$
and $\applyfunc{A}{t} = A$ and so $\applyfunc{x}{t} = x$ for all $x \in A$. 

\begin{lemma}
\label{app:lem:SIsets}
	Let $\semiautomat{A}{Q}{\Sigma}{\delta}$ be an SI-automaton,
	$A,B,C \subseteq Q$ be $k$-sets, and let $u \in \Sigma^*$
	be a shortest word such that $\applytrans{\delta}{A}{u} = C$ or $\applytrans{\delta}{B}{u} = C$.
	If $A \ne B$, then $\applytrans{\delta}{A}{u} \ne \applytrans{\delta}{B}{u}$.
\end{lemma}
\begin{proof}
	Suppose $\applytrans{\delta}{A}{u} = \applytrans{\delta}{B}{u}$.
	Next, we prove $A = B$
	by induction on $m = \sum_{a \in \Sigma_I} |u|_a$.
	If $m = 0$,
	then $u \in \Sigma_G^*$.
	So, as $u$ permutes the states, 
	$\applytrans{\delta}{A}{u} = \applytrans{\delta}{B}{u}$
	implies $A = B$.
	Otherwise, write $u = vaw$ with $v \in \Sigma^*$, $a \in \Sigma_I$, and $w \in\Sigma_G^*$.
	Then $\applytrans{\delta}{A}{va} = \applytrans{\delta}{B}{va}$ since~$w$ permutes the states.
	As $u$ has shortest length,
	we must have $\applytrans{\delta}{\applytrans{\delta}{A}{v}}{a} \ne \applytrans{\delta}{A}{v}$
	and $\applytrans{\delta}{\applytrans{\delta}{B}{v}}{a} \ne \applytrans{\delta}{B}{v}$,
	as otherwise $vu$ would give a shorter word mapping at least one of $A$ or $B$ onto $C$.
	Since $|C| \le |\delta(A, v)| \le |A|$ and $|C| \le |\delta(B, v)| \le |B|$ by \cref{eqn:subset-mapped-to-smaller-set},
	and $|A| = |B| = |C|$
	by assumption,
	we have $|\applytrans{\delta}{A}{v}| = |\applytrans{\delta}{B}{v}|$.
	So, by \cref{app:lem:SI} applied to $\fword{\delta}{a}$,
	we get $\applytrans{\delta}{A}{v} = \applytrans{\delta}{B}{v}$. 
	Since~$u$ was chosen minimal such that $\delta(A, u) = C$
	or $\delta(B, u) = C$, also~$v$ has shortest length
	such that $\delta(A, v) = C'$ or $\delta(B, v) = C'$
	for $C' = \delta(A, v) = \delta(B, v)$.	
	By induction hypothesis, we conclude that $A = B$.
\end{proof}

\begin{remark}
	The assumption that the sets $A, B, C$
	have the same size in \cref{app:lem:SIsets} 
	can not be dropped.
	For example, consider the SI-automaton
	with states $\{1,2,3\}$ and two simple idempotent letters $a$
	and $b$ that map $1$ to~$2$
	and $2$ to $3$, respectively.
	Then the word $ab$ maps the two distinct sets $\{1,2\}$
	and $\{2,3\}$ to $\{3\}$.	
\end{remark}

\begin{corollary} \label{app:cor:SIsets}
	Let $\semiautomat{A}{Q}{\Sigma}{\delta}$ be an SI-automaton
	and $p,q \in Q$ be distinct states.
	Then for a shortest word $u \in \Sigma^*$
	with $\delta(p, u) = r$ or $\delta(q, u) = r$ for some $r \in Q$
	we have $\delta(p, u) \ne \delta(q, u)$.	
\end{corollary}
\begin{proof}
	This is \cref{app:lem:SIsets} for the special case $k = 1$.	
\end{proof}

\begin{proposition}
\label{app:prop:pairof2sets}
	Let $\semiautomat{A}{Q}{\Sigma}{\delta}$ be an SI-automaton
	and $A, B \subseteq Q$ be two distinct $2$-sets.
	Then for a shortest word $u \in \Sigma^*$
	with $|\applytrans{\delta}{A}{u}| = 1$
	or $|\applytrans{\delta}{A}{u}| = 1$
	we have $|\applytrans{\delta}{A}{u}| \ge 2$
	or $|\applytrans{\delta}{A}{u}| \ge 2$.	
\end{proposition} 
\begin{proof}
	Let $q \in Q$ be the unique state with $\applytrans{\delta}{q}{a}\ne q$.
	As $u$ has shortest length mapping at least one to a singleton
	subset, we can write $u = va$ with $v \in \Sigma^*$
	and $a \in \Sigma_I$
	where $v$ has minimal length such that
	$\applytrans{\delta}{A}{v} = \{q, \applytrans{\delta}{q}{a}\}$
	or $\applytrans{\delta}{B}{v} = \{q, \applytrans{\delta}{q}{a}\}$.	
	Now, suppose $|\applytrans{\delta}{A}{u}| = |\applytrans{\delta}{B}{u}| = 1$.
	By the minimal choice of $u$,
	we must have $|\applytrans{\delta}{A}{v}| = |\applytrans{\delta}{B}{v}| = 2$.
	Since both sets are mapped to a singleton, we find
	$\applytrans{\delta}{A}{v} = \applytrans{\delta}{B}{v} = \{q, \applytrans{\delta}{q}{a}\}$.
	Then we have $|A| = |B| = |\delta(A, v)|$.
	Therefore, by \cref{app:lem:SIsets}, since~$v$ has shortest length mapping
	at least one to $\{q, \delta(q, a)\}$, we conclude $A = B$.	
	So, with the assumption $A \ne B$,
	we find that at least one of 
	both subsets is mapped to a non-singleton set.
\end{proof}

\begin{corollary}\label{syncmax:cor:SI-2-sets}
	In a synchronizing SI-automaton every two distinct $2$-sets
	are distinguishable.
\end{corollary}
\begin{proof}
 By Proposition~\ref{app:prop:pairof2sets}, a shortest
 word mapping at least one of two distinct $2$-sets
 to a singleton set must map the other set to a non-singleton
 set. Hence, it distinguishes both $2$-sets.
\end{proof}

Recall that \cref{sync_max:ex:2-sets-not-all-dist} gave synchronizing SI-automata
that are not completely distinguishable.

\subsection{Strongly Connected SI-Automata}
\label{sec:SI:SCC}


Here, in \cref{sync_max:thm:SI-automata},
we show that for strongly connected SI-automata,
synchronizability, complete reachability,
and distinguishability of pairs of $2$-sets
are equivalent properties. As a consequence, stated in \cref{sync_max:cor:SIcr},
completely reachable SI-automata are precisely those that
are synchronizing and strongly connected.
\Cref{cor:SI_CR} and \cref{prelim:classes:SI:without-perm-letters-SCC} 
link strong connectivity and strong trap-connectivity
to complete reachability and absence of non-trivial congruences
in SI-automata without permutational letters.
\cref{syncmax:SI-aut:cor:single-idempotent} states various equivalences for SI-automata with a single non-permutational letter,
and \cref{sync_max:prop:binary-SI} for binary SI-automata.

\begin{theorem}
	\label{sync_max:thm:SI-automata}
	For a strongly connected $n$-state SI-automaton $\semiautomat{A}{Q}{\Sigma}{\delta}$ with $n \ge 3$ the following are equivalent.
	\begin{enumerate}
		\item \label{sync_max:thm:SI-automata:1} $\mathcal A$ is synchronizing.
		
		\item \label{sync_max:thm:SI-automata:2} $\mathcal A$ is completely reachable. 
		
		\item \label{sync_max:thm:SI-automata:4} Every automaton for $\Syn(\mathcal A)$ needs $\binom{n}{2} + 1$ states.
		
		\item \label{sync_max:thm:pairs-2-sets} All pairs of distinct $2$-sets of states are distinguishable.
	\end{enumerate}
\end{theorem}
\begin{proof}
	$\ref{sync_max:thm:SI-automata:1} \Rightarrow \ref{sync_max:thm:SI-automata:2}$:
	Suppose $\mathcal A$ is synchronizing. 
	Let $I(\mathcal A) = (Q, E)$ be the Rystsov graph of $\mathcal A$. 
	By \cref{sync_max:thm:rystsov_SI}
	this graph is strongly connected.
	By \cref{sync_max:prop:Rystsov_graph_vs_SI_graph}
	it is a subgraph of $I(\TM{\mathcal A})$.
	From Theorem~\ref{sync_max:thm:IS_SCC}
	it follows that~$\TM{\mathcal A}$, and so $\mathcal A$,
	is completely reachable.\footnote{For an alternative proof:
			Let $a_1, \ldots, a_k$ be the simple idempotent letters where we have $k \ge 1$ since $\mathcal A$ is synchronizing.
			By assumption, for each $q \in Q$
			there exists $t \in \TM{\mathcal A}$
			such that $t(Q) = \{q\}$.
			We can write $t = g_0 a_{i_1} g_1 a_{i_2} \cdots g_{m-1} a_{i_m} g_m$
			for permutations $g_0, g_1, \ldots, g_m \in \TM{\mathcal A}$
			and numbers $1 \le i_1, i_2, \ldots, i_m \le k$
			where we write $a_i$ for~$\fword{\delta}{a_i}$ here as a shorthand.
			Since 
			\[
			t = a_{i_1}^{g_0} a_{i_2}^{g_0 g_1} \cdots  a_{i_m}^{g_0 g_1 \cdots g_m} g_0 g_1 \cdots g_m,
			\]
			we conclude that $a_{i_1}^{g_0} a_{i_2}^{g_0 g_1} \cdots  a_{i_m}^{g_0 g_1 \cdots g_m}$
			already maps $Q$ to $\{q\}$.
			By considering the transformations $a_i^g$ for a permutation $g \in \TM{\mathcal A}$
			as letters in a new automaton $\mathcal B$ (Automata constructed
					by introducing new symbols for existing words or transformations were also considered in~\cite[p.~270]{AnaVolGus2013}.
					We will also define automata in this way in \cref{sync_max:chap:slowly}, for example, in the proof of \cref{syn_max:thm:Mn}.), we find 
			that $\mathcal B$ must be synchronizing and, since each state is synchronizing,
			strongly connected.
			Each $a_i^g$ induces a letter of simple idempotent,
			and so it induces an edge in $\Gamma_1(\mathcal A)$.
			Now, the automaton graph of~$\mathcal B$
			with its self-loops omitted is a subgraph of $\Gamma_1(\mathcal A)$
			since the unique transition of a simple idempotent that is not a self-loop
			correspond to an edge of $\Gamma_1(\mathcal A)$ (this does not hold in general, since some transitions for a letter might not induce edges in $\Gamma_1(\mathcal A)$).
			So $\Gamma_1(\mathcal A)$ is strongly connected.
			By \cref{sync_max:thm:char_compl_reach}
			we find that~$\mathcal B$, and so~$\mathcal A$,
			is completely reachable.
	}
	
	$\ref{sync_max:thm:SI-automata:2} \Rightarrow \ref{sync_max:thm:SI-automata:4}$:
	Since a completely reachable automaton is synchronizing, from \cref{syncmax:cor:SI-2-sets}
	we get that distinct pairs of $2$-sets are distinguishable
	(alternatively, it also follows 
	from \cref{sync_max:thm:IS_SCC}
	as in the proof of the previous implication).
	Furthermore, all $2$-sets are reachable
	and so $\{ \{p,q\} \mid p,q \in Q, p \ne q \} \cup \{\{s\}\}$ for some $s \in Q$
	yields a family of pairwise distinguishable
	and reachable subsets. By \cref{sync_max:thm:lower_bound_SynA}, 
	it follows that $\Syn(\mathcal A)$
	has at least $\binom{n}{2} + 1$ states.
	
	$\ref{sync_max:thm:SI-automata:4} \Rightarrow \ref{sync_max:thm:pairs-2-sets}$:
	As $n \ge 3$
	the minimal automaton for $\Syn(\mathcal A)$ has at least
	five states. 
	So, $\Syn(\mathcal A) \ne \emptyset$, i.e., $\mathcal A$ is synchronizing,
	and from \cref{syncmax:cor:SI-2-sets}
	we get that distinct pairs of $2$-sets are distinguishable.
	
	$\ref{sync_max:thm:pairs-2-sets} \Rightarrow \ref{sync_max:thm:SI-automata:1}$:
	As $\mathcal A$ is strongly connected and $n \ge 3$, by 
	Proposition~\ref{sync_max:prop:2-sets-dist-SCC-implies-sync}
	the automaton $\mathcal A$ is synchronizing. 
\end{proof}

\Cref{sync_max:fig:sync-two-sets-dist} 
gives examples of strongly connected and completely
reachable SI-automata that are not completely distinguishable (and so not sync-maximal).

\Cref{sync_max:thm:SI-automata}
fails for synchronizing SI-automata
that are not strongly connected since they are not completely reachable,
as, for example, the SI-automaton from \cref{prelim:fig:binarySI-sync-trap}.

\begin{corollary}\label{sync_max:cor:SIcr}
	An SI-automaton is strongly connected
	and synchronizing if and only if it is completely reachable.
\end{corollary}
\begin{proof}
	By \crefitem{sync_max:thm:SI-automata}{sync_max:thm:SI-automata:2},
	a strongly connected and synchronizing
	SI-automaton is completely reachable.
	\Cref{prelim:lem:compl_reach_implies_SCC_sync} yields the other implication.
\end{proof}

Strong connectivity and synchronizability do not imply complete reachability
even for automata where every letter has defect at most one,
see \cref{syncmax:fig:SCC-sync-not-cr}.

\begin{figure}[htb]
	\centering
\begin{tikzpicture}[->, >=stealth, node distance=3cm, on grid, auto,
	state/.style={circle, draw, minimum size=6mm, inner sep=1pt}
	]
	\node[state] (q1) {};
	\node[state] (q2) [right=of q1] {};
	\node[state] (q3) [right=of q2] {};
	
	\path
	(q1) edge [bend left=15] node {$a, b$} (q2)
	(q2) edge [bend left=15] node {$a$} (q3)
	(q3) edge [bend left=15] node {$a,b$} (q2)
	(q2) edge [bend left=15] node {$b$} (q1);
\end{tikzpicture}
	\caption{Strongly connected and synchronizing automaton that is not completely reachable.}
	\label{syncmax:fig:SCC-sync-not-cr}
\end{figure}

\begin{corollary} \label{syncmax:cor:SI-transitive-grp-of-units}
	Let $\semiautomat{A}{Q}{\Sigma}{\delta}$
	be an SI-automaton such that the group of units of $\TM{\mathcal A}$ is transitive.
	Then the following are equivalent.
	
	\begin{enumerate}
		\item $\mathcal A$ is synchronizing.
		
		\item $\mathcal A$ is completely reachable.
		
		\item $\mathcal A$ is completely distinguishable.
		\end{enumerate}
\end{corollary}

Note that we also have the equivalences from \cref{cor:defectatmostone-weaker} for SI-automata
where the group of units of its transformation monoid is transitive.

\begin{corollary}
	\label{cor:SI_CR}
	Let $\semiautomat{A}{Q}{\Sigma}{\delta}$ be an SI-automaton without permutational letters.
	Then $\mathcal A$ is completely reachable
	if and only if it is strongly connected.
\end{corollary}
\begin{proof}
	Suppose $\semiautomat{A}{Q}{\Sigma}{\delta}$ is a strongly connected
	SI-automaton without permutational letters.	
	Let $A \subseteq Q$ be a maximal subset such that $|\delta(A, u)| = 1$
	for some $u \in \Sigma^*$.
	Assume $|A| < |Q|$.
	As $\mathcal A$ is strongly connected,
	there exists a transition labeled by some letter $a \in \Sigma$
	starting from some state in $Q \setminus A$
	and ending in $A$, i.e. $\delta(p, a) \in A$
	for some $p \in Q \setminus A$.
	As~$a$ is a simple idempotent,
	we have $\delta(A \cup \{p\}, a) = A$.
	Then $|\delta(A \cup \{p\}, au)| = 1$, contradicting
	the maximal choice of~$A$.
	So, we must have $A = Q$ and $\mathcal A$
	is synchronizing.		
	Then the claim follows from \cref{sync_max:cor:SIcr}.\footnote{For an alternative argument,
	for such an automaton, the automaton graph with its self-loops omitted
	equals its Rystsov graph $I(\mathcal A)$. Thus it is strongly connected
	and so $\mathcal A$ is completely reachable by \cref{sync_max:thm:IS_SCC}.}
\end{proof}

\begin{proposition}  \label{sync_max:prop:sufficient-primitive-SI}
	Let $\semiautomat{A}{Q}{\Sigma}{\delta}$ be an SI-automaton with at least three states.
	Then $\mathcal A$ does not admit some non-trivial $\TS{\mathcal A}$-congruence and has some simple idempotent letter
	if and only if it is strongly trap-connected, or it is
	strongly connected and synchronizing.
\end{proposition}
\begin{proof}
	First, assume $\mathcal A$ does not admit some non-trivial $\TS{\mathcal A}$-congruence and 
	has some simple idempotent letter.
	By~\cite[Cor.~2.3]{Volkov24} an automaton not admitting non-trivial congruences
	and having a letter inducing a one-point singular is synchronizing.
	Thus, $\mathcal A$ is synchronizing.
	By \cref{prelim:cor:simpleaut}, 
	it is strongly connected or strongly trap-connected.	
	
	Conversely, first suppose~$\mathcal A$ is strongly trap-connected with trap state $t$.
Let $\sim$ be some $\TS{\mathcal A}$-congruence
and suppose $p \sim q$ for distinct states $p,q \in Q$.
Consider a shortest word $u$
mapping $p$ or $q$ to~$t$.
By \cref{app:cor:SIsets} it
maps only one of both states to~$t$
and the other state remains in~$Q \setminus \{t\}$.
Without loss of generality,
let $\delta(p, u) = t$
and $\delta(q, u) \in Q \setminus \{t\}$.
Let $r \in Q$ be some state with $r \ne t$.
Since~$Q \setminus \{t\}$ is strongly connected,
we have $\delta(q, uv) = r$
for some $v \in \Sigma^*$.
With $\delta(p, u) \sim \delta(q, u)$, this implies $t \sim r$.
Since~$r$ was chosen arbitrarily in~$Q \setminus \{t\}$,
we find that every state is congruent to~$t$
and thus~$\sim$ is the universal relation.

Next, suppose~$\mathcal A$
is strongly connected and synchronizing.
By \cref{sync_max:thm:rystsov_SI}, $I(\mathcal A)$
is strongly connected. Since $I(\mathcal A)$ is a subgraph of $I(\TS{\mathcal A})$
by \crefitem{sync_max:prop:Rystsov_graph_vs_SI_graph}{sync_max:prop:item:IAsubgraphIS},
from \cref{thm:primitive} we get that $\mathcal A$
admits no non-trivial $\TS{\mathcal A}$-congruence. 
\end{proof}


	SI-automata without simple idempotent letters
	whose group of units is primitive do not admit non-trivial congruences
	and are not synchronizing (and so also not strongly trap-connected).
	
	
	\begin{corollary} \label{prelim:classes:SI:without-perm-letters-SCC}
		Let $\semiautomat{A}{Q}{\Sigma}{\delta}$ be an SI-automaton with at least three states without permutational letters.
		Then~$\mathcal A$ does not admit some non-trivial $\TS{\mathcal A}$-congruence
		if and only if it is strongly trap-connected or strongly connected.	
	\end{corollary}
	\begin{proof}
		If~$\mathcal A$
		does not admit any non-trivial $\TS{\mathcal A}$-congruence,
		then, by \cref{prelim:cor:simpleaut}, 
		it is strongly connected or strongly trap-connected.	
		The converse implication follows from \cref{sync_max:prop:sufficient-primitive-SI} 
		and the fact that if $\mathcal A$ is strongly connected, then it is synchronizing
		by \cref{cor:SI_CR}.
		%	% Alter, direkter Beweis ohne die \cref{sync_max:prop:sufficient-primitive-SI} 
		%		First, suppose~$\mathcal A$ is strongly connected or strongly trap-connected.
		%		Let $A \subseteq Q$ be a proper subset with $|A| \ge 2$.
		%		There exists an edge connecting a state in~$A$
		%		to some state state outside of~$A$, i.e., $\delta(q, a) \notin A$
		%		for some $q \in A$ by the following argument: if $\mathcal A$ is strongly connected,
		%		this follows from \cref{prelim:lem:partition_SCC},
		%		if $\mathcal A$ is strongly trap-connected with trap state~$t$,
		%		then for $A \ne Q \setminus \{t\}$
		%		this follows from \cref{prelim:lem:partition_SCC}
		%		and for $A = Q \setminus\{t\}$
		%		we can choose $\delta(q, a) = t$.
		%		Let $p \in A$ with $p \ne q$.
		%		Then $p, q \in A$ but $\delta(p, a) = p \in A$
		%		and $\delta(q, a) \notin A$ since~$a$ is a simple idempotent.
		%		Thus, $A$ cannot be the congruence class of some $\TS{\mathcal A}$-congruence.
		%		Since~$A$ was arbitrarily chosen, this implies
		%		that $\mathcal A$ only admits the trivial $\TS{\mathcal A}$-congruences.
	\end{proof}
	

\begin{proposition} \label{syncmax:prop:SI-grp-units-transitive}
	Let $\semiautomat{A}{Q}{\Sigma}{\delta}$
	be an SI-automaton with at least three states and 
	a single simple idempotent letter.
	Then the following implies that the group of units $G$ of $\TM{\mathcal A}$ is transitive:
	
	\begin{enumerate}
	\item \label{syncmax:prop:SI-grp-units-transitive:1} $\mathcal A$ is strongly connected.
	
	\item \label{syncmax:prop:SI-grp-units-transitive:2} Every $k$-set is distinguishable from every distinct $l$-set of states
	for some $2 \le k, l \le |Q|$ with $\min\{k,l\} < |Q|$ and $\max\{k, l\} > 2$.
	\end{enumerate}
\end{proposition}
\begin{proof}
\ref{syncmax:prop:SI-grp-units-transitive:1}:	Suppose $G$ is not transitive and let $\Orbit{p}{G}$
	and $\Orbit{q}{G}$ be two distinct orbits for $p,q \in Q$.
	If the single idempotent letter connects two states
	in one of these orbits, then no state from one orbit
	is reachable from a state from the other orbit
	and so $\mathcal A$ is not strongly connected.
	Otherwise, if, without loss of generality,
	this letter connects a state from $\Orbit{p}{G}$
	with a state from $\Orbit{q}{G}$,
	then no state from  $\Orbit{p}{G}$
	is reachable from a state from  $\Orbit{q}{G}$
	and so~$\mathcal A$ is not strongly connected.
	
\ref{syncmax:prop:SI-grp-units-transitive:2}: Without loss of generality, suppose $k \le l$. 
Let $a \in \Sigma$ be the simple idempotent letter
and $\delta(p, a) = q$ with $p \ne q$.
From \cref{prelim:autsyncw:prop:two-sets-min-rank} 
we get that~$\mathcal A$ is synchronizing.
The orbits $\Orbit{p}{G}$ and $\Orbit{q}{G}$ are the only orbits of~$G$,
for if there exists another orbit, then it is mapped to itself by every letter
and no letter maps some state into it, which
implies that~$\mathcal A$ cannot be synchronizing.
For $\Orbit{p}{G} \ne \Orbit{q}{G}$,
there is no letter mapping two distinct states in~$\Orbit{q}{G}$
to one state and so $\Orbit{q}{G} = \{q\}$
and~$q$ is a trap state.
We can pick distinct $s, t \in A$ as $|Q| \ge 3$.
Since $l > 2$ and $2 \le k < |Q|$
there exists a $k$-set $A$ with $s,t \in A$
and an $l$-set $B$ with $q, s,t \in B$ such that $A \ne B$.
The only way to map two states to one state is to use the letter~$a$,
which yields that a word mapping two states to one state maps them both to~$q$.
Thus, $|\delta(A, u)| = 1$
if and only if $\delta(A, u) = \{q\}$
if and only if $\delta(B, u) = \{q\}$
if and only if $|\delta(B, u)| = 1$,
which yields that~$A$ and~$B$ are indistinguishable.
So, $G$ must be transitive.
\end{proof}

\begin{corollary}\label{syncmax:SI-aut:cor:single-idempotent}
	Let $\semiautomat{A}{Q}{\Sigma}{\delta}$
	be an SI-automaton with at least three states and a single simple idempotent letter.
	Then the following are equivalent.
	\begin{enumerate}
		\item \label{syncmax:SI-aut:cor:single-idempotent:1} $\mathcal A$ is strongly connected and synchronizing. 
		
		\item \label{syncmax:SI-aut:cor:single-idempotent:2} $\mathcal A$ is completely reachable.
		
		\item \label{syncmax:SI-aut:cor:single-idempotent:3} $\mathcal A$ is sync-maximal.
		
		\item \label{syncmax:SI-aut:cor:single-idempotent:4} $\mathcal A$ is completely distinguishable.		

		\item \label{syncmax:SI-aut:cor:single-idempotent:5} Every $k$-set is distinguishable from every distinct $l$-set of states for some $2 \le k, l \le |Q|$ with $\min\{k,l\} < |Q|$ and $\max\{k, l\} > 2$.
		
		\item \label{syncmax:SI-aut:cor:single-idempotent:6} $\mathcal A$ is strongly connected and all pairs of distinct $2$-sets are distinguishable.
	\end{enumerate}
\end{corollary}
\begin{proof}
	$\ref{syncmax:SI-aut:cor:single-idempotent:1} \Rightarrow \ref{syncmax:SI-aut:cor:single-idempotent:2}$:
	This follows from \cref{sync_max:cor:SIcr}, 
	
	$\ref{syncmax:SI-aut:cor:single-idempotent:2} \Rightarrow \ref{syncmax:SI-aut:cor:single-idempotent:3}$:	
	As it has only one non-permutational letter, by \cref{sync_max:cor:single-non-permutation}, 
	we find that the group of units of of $\TM{\mathcal A}$  is transitive.\footnote{Alternatively,
		this also follows from \cref{syncmax:prop:SI-grp-units-transitive},
		or with the following argument: for a completely reachable SI-automaton $\mathcal A$, by \cref{prelim:lem:compl_reach_implies_SCC_sync} and \cref{sync_max:thm:rystsov_SI},~$I(\mathcal A)$ 
		is strongly connected. In case $\mathcal A$ has only a single simple idempotent mapping~$s$ to~$t$,~$I(\mathcal A)$ equals
		the directed orbital graph of $(s,t)$ w.r.t.~the group of units $G$ of $\TM{\mathcal A}$.
		Then, strong connectivity implies
		that each $q \in Q$ has an outgoing edge $(g(s), g(t))$, $g \in G$.
		So, $q = g(s)$ for some $g \in G$ and so $G$ is transitive.}
	So, since $|Q| \ge 3$,
	$\mathcal A$ is completely distinguishable by \cref{cor:defectatmostone-weaker}
	and thus it is sync-maximal by \cref{sync_max:cor:SCC}.\footnote{For an alternative proof of complete distinguishability:
		Let $p, q \in Q$
		with $\excl{a} = \{p\}$ and $\dupl{a} = \{q\}$
		for the unique simple idempotent letter $a \in \Sigma$.
		Consider the automaton~$\mathcal B$
		with letters $a^g$
		where $a^g$ acts as the conjugate of the action of the letter~$a$
		for some $g$ from the group of units of $\TM{\mathcal A}$.
		Since~$\mathcal A$ is completely reachable, as in the proof of \cref{sync_max:thm:SI-automata}
		presented in the footnote, we find that this automaton is strongly connected.
		Since $a^g$ maps $g(p)$ to $g(q)$ by \cref{prelim:conj},
		the automaton graph of $\mathcal B$ equals the orbital digraph
		of $(p,q)$ of the group of units of $\TM{\mathcal A}$
		and it also equals $I(\mathcal A)$.
		In this orbital digraph, by strong connectivity,
		every state has an outgoing transition, which
		implies that the group of units must be transitive
		since then there must exist a permutation mapping~$p$
		to every other state.
		Thus it is vertex-transitive and since $p \ne q$ it has no self-loops, by \cref{prelim:prop:vertex-transitive-SCC-32-connected},
		we find that the orbital digraph is strongly $\frac{3}{2}$-connected.
		Since $I(\mathcal A)$ is a subgraph of $I(\TM{\mathcal A})$
		by \cref{sync_max:prop:Rystsov_graph_vs_SI_graph},
		from \cref{sync_max:thm:IS_32SCC}, 
		we find that~$\mathcal A$
		is completely distinguishable.}	
	
	$\ref{syncmax:SI-aut:cor:single-idempotent:3} \Rightarrow \ref{syncmax:SI-aut:cor:single-idempotent:4}$: By \cref{sync_max:thm:char_sync_max_aut}, $\mathcal A$ is completely distinguishable.
	
	$\ref{syncmax:SI-aut:cor:single-idempotent:4} \Rightarrow \ref{syncmax:SI-aut:cor:single-idempotent:5}$: This is obvious.
	
	$\ref{syncmax:SI-aut:cor:single-idempotent:5} \Rightarrow \ref{syncmax:SI-aut:cor:single-idempotent:6}$:  
	By \cref{syncmax:prop:SI-grp-units-transitive}, the group of units of $\TM{\mathcal A}$
	is transitive.
	Since every state can be mapped to any other state
	by a permutational word, $\mathcal A$
	is strongly connected.
	By \cref{sync_max:thm:SI-automata} we find
	that all pairs of distinct $2$-sets are distinguishable.	
	
	$\ref{syncmax:SI-aut:cor:single-idempotent:6} \Rightarrow \ref{syncmax:SI-aut:cor:single-idempotent:1}$: 
	This follows from \cref{sync_max:prop:2-sets-dist-SCC-implies-sync}.
\end{proof}

Next, we state equivalences for binary SI-automata
derived from \cref{syncmax:SI-aut:cor:single-idempotent}.

\begin{corollary} \label{sync_max:prop:binary-SI}
	Let  $\semiautomat{A}{Q}{\{a,b\}}{\delta}$ be a binary SI-automaton with at least three states.
	Then the following are equivalent.
	
	\begin{enumerate}
		\item \label{sync_max:cor:binary-SI:1} $\mathcal A$ is strongly connected and synchronizing.
		
		\item \label{sync_max:cor:binary-SI:2} $\mathcal A$ is completely reachable.
		
		\item \label{sync_max:cor:binary-SI:3} $\mathcal A$ is sync-maximal.
		
		\item \label{sync_max:cor:binary-SI:4a} $\mathcal A$ is completely distinguishable.
		
		\item \label{sync_max:cor:binary-SI:4b} Every $k$-set is distinguishable from every distinct $l$-set of states for some $2 \le k, l \le |Q|$ with $\min\{k,l\} < |Q|$ and $\max\{k, l\} > 2$.
			
		\item \label{sync_max:cor:binary-SI:5} $\mathcal A$ is strongly connected and all pairs of distinct $2$-sets are distinguishable.
		
		\item \label{sync_max:cor:item:form} In $\mathcal A$, one letter, say $b \in \Sigma$, cyclically permutes
		the states and the other letter $a \in \Sigma$
		is a simple idempotent 
		such that $\applytrans{\delta}{q}{a} = \applytrans{\delta}{q}{b^d a}$
		for some $q \in Q$ and $0 < d < n$ coprime to~$n$ (see \cref{prelim:fig:binarySI-sync-SCC}).
	\end{enumerate}
\end{corollary} 
\begin{proof}
	It is easily seen that the only synchronizing binary SI-automata
	without permutational letters and $|Q| \ge 3$
	have three states and fulfill neither of the
	mentioned properties. Thus, we can assume~$\mathcal A$
	has a single simple idempotent letter.
	Then, by \cref{syncmax:SI-aut:cor:single-idempotent},
	we obtain the equivalences
	between~\cref{sync_max:cor:binary-SI:1,sync_max:cor:binary-SI:2,sync_max:cor:binary-SI:3,sync_max:cor:binary-SI:4a,sync_max:cor:binary-SI:4b,sync_max:cor:binary-SI:5}.
	
	A synchronizing binary SI-automaton has
	only the two forms from \cref{prop:structure_simp_idem_binary},
	where only the type of automaton as claimed 
	is strongly connected and synchronizing. 
	This gives the equivalence between \cref{sync_max:cor:binary-SI:1,sync_max:cor:item:form}.
\end{proof}

In the binary automaton from \cref{prelim:fig:binarySI-sync-trap},
pairs of distinct $2$-set are distinguishable, 
but it is neither strongly connected nor completely distinguishable.

\section{Circular Automata}
\label{sync_max:sec:circular}

%!TEX root = ./../../thesis-trier.tex

\IN{circular automaton}
An automaton $\semiautomat{A}{Q}{\Sigma}{\delta}$ is \emph{circular}
if there exists a letter $b \in \Sigma$ 
that cyclically permutes the states, i.e., $Q = \{ \delta(q, b^i) \mid i \ge 0 \}$
and $\delta(q, b^{|Q|}) = q$ for some state $q \in Q$ (and hence all states $q \in Q$).
The automata from the \Cerny family (see Figure~\ref{intro:fig:cerny}) are circular.
Dubuc~\cite{Dubuc98} confirmed the \Cerny conjecture for circular automata.

In \cref{syncmax:classes:thm:circular}, we first give a condition implying complete reachability,
complete distinguishability, and bounds on the lengths of shortest words reaching non-empty subsets.
This result was stated by Don~\cite[Prop.~15]{Don16} for the special case of a single
letter of defect one and without concluding complete distinguishability.
We then derive several consequences and give examples.

\begin{theorem} \label{syncmax:classes:thm:circular} \label{sync_max:prop:don}
	Let $\semiautomat{A}{Q}{\Sigma}{\delta}$ be an $n$-state automaton with $\Sigma = \{a_1, \ldots, a_k, b\}$
	where $b \in \Sigma$ cyclically permutes the states and
	$a_1, \ldots, a_k \in \Sigma$ are letters of defect one
	with $0 < d_1, \ldots, d_k < n$	such that
	\begin{enumerate}
	\item for all $i \in \{1,\ldots,k\}$ we have $\dupl{a_i} = \delta(\excl{a_i}, b^{d_i})$ and
	
	\item the greatest common divisor of $d_1, \ldots, d_k$ and $n$ is one.
	\end{enumerate}
	Then $\mathcal A$ is completely distinguishable,
	every non-empty subset of size~$r$ is reachable by a word of length $n(n - r)$,
	and a shortest synchronizing word has length at most $(n - 1)^2$.
\end{theorem}
\begin{proof}
	Consider the digraph $\Gamma_{1,n}(\mathcal A)$ defined on page~\pageref{def:Gamma1i-graph}.
	First, we show that the group of units~$G$ of $\TM{\mathcal A}$ is contained in the 
	automorphism group of $\Gamma_{1,n}(\mathcal A)$.
	Since~$G$ is generated 
	by the action of the letter~$b$, this follows if for every edge $(p,q)$ in $\Gamma_{1,n}(\mathcal A)$
	also the edge $(\delta(p, b), \delta(q, b))$ is in $\Gamma_{1,n}(\mathcal A)$.	
	Let $(p, q)$ be an edge in $\Gamma_{1,n}(\mathcal A)$.
	Then it is induced by a word of the form $a_i b^{j}$
	for some $i \in \{1,\ldots,k\}$ and $j \in \{0,1,\ldots,n-1\}$.
	For the word $u = a_i b^{(j+1) \bmod n}$
	we have $\excl{u} = \{ \delta(p, b)\}$
	and $\dupl{u} = \{ \delta(q, b) \}$.
	Thus, $u$ induces the edge $(\delta(p, b), \delta(q, b))$ in $\Gamma_{1,n}(\mathcal A)$.
		
	Next, assume $\Gamma_{1,n}(\mathcal A)$ is not strongly connected.
	An automorphism maps a path to another path.
	Thus, since the elements from~$G$ act as automorphisms, 
	the weakly connected components form the congruence classes of
	a $G$-congruence on~$Q$.
	Then, by \cref{lem:congruence_on_cyclic_group} there exists a divisor $d$ of~$n$ such that
	two states $p,q \in Q$ are in a common weakly connected component
	if and only if $\delta(p, b^j) = q$
	for some multiple~$0 \le j < n$ of~$d$.
	Note that since~$d$ divides~$n$, the numbers
	$j$ and $n - j$ are either both divisible by~$d$
	or neither is divisible by~$d$.	
	For $i \in \{1,\ldots,k\}$,
	let $\excl{a_i} = \{p_i\}$
	and $\dupl{a_i} = \{q_i\}$.
	Since $(p_i, q_i)$ is an edge in $\Gamma_{1,n}(\mathcal A)$,
	the states~$p_i$ and~$q_i$ are in the same weakly connected component.
	Thus $\delta(p_i, b^j) = q_i$ for some multiple $0 \le j < n$ of~$d$.
	Since $\delta(p_i, b^{d_i}) = q_i$,
	we find $j = d_i$.
	This yields that~$d$ is a common divisor
	of $d_1, \ldots, d_k$ and~$n$.
	By assumption, we find $d = 1$. This implies that all states
	are congruent, i.e., $\Gamma_{1,n}(\mathcal A)$ is weakly connected.
	Since the group $G$ is transitive, $\Gamma_{1,n}(\mathcal A)$ is vertex transitive.
	Thus, by \cref{prelim:thm:vertex_transitive_finite_graph}, it is strongly connected.
	Then, by \cref{syncmax:thm:Gamma1-SCC-subgraph}, 
	$\mathcal A$ is completely reachable,
	each non-empty subset of size $r$ is reachable by a word of length $n(n - r)$,
	and a shortest synchronizing word has length at most $1 + n(n-2) = (n-1)^2$.
	By \cref{thm:defectatmostone}, we find that~$\mathcal A$ is completely distinguishable.
\end{proof}

The following corollary was stated in~\cite[Thm.~13]{HoffmannCircularExtended} in the special case $k = 1$.

\begin{corollary}  \label{sync_max:thm:circular_sync_max} \label{sync_max:cor:circular}
	Let $\semiautomat{A}{Q}{\Sigma}{\delta}$ be an $n$-state automaton where $b \in \Sigma$ cyclically permutes the states.
	Suppose $a_1, \ldots, a_k \in \Sigma$ are letters of defect one
	and $0 < d_1, \ldots, d_k < n$
	are such that for all $i \in \{1,\ldots,k\}$ we have $\delta(q_i, a_i) = \delta(q_i, b^{d_i} a_i)$
	for some $q_i \in Q$.
	If the greatest common divisor of $d_1, \ldots, d_k$ and $n$
	is one, then $\mathcal A$ is completely reachable and completely distinguishable. % (and so sync-maximal).	
\end{corollary}
\begin{proof}
	Since $\TM{\mathcal A}$ contains the transitive permutation group $G = \langle \{ \fword{\delta}{u} \} \rangle_G$,
	by \cref{sync_words:lem:idempotent}, for each $i \in \{1,\ldots,k\}$
	there exist simple idempotent words $u_i$
	mapping $q_i$ to $\delta(q_i, b^{d_i})$.
	We have $\excl{u_i} = \{q_i\}$
	and $\dupl{u_i} = \{ \delta(q_i, b^{d_i}) \}$.
	Considering a new automaton with state set~$Q$ with the letter~$b$
	acting in the same way and where we consider the words $u_1, \ldots, u_k$ as letters,
	from \cref{syncmax:classes:thm:circular} 
	we get that
	this automaton is completely reachable and completely distinguishable.
	This also yields that~$\mathcal A$ is completely reachable and completely distinguishable.
\end{proof}

We will apply \cref{sync_max:thm:circular_sync_max} 
in \cref{sync_max:sec:slowly_sync} (see \cref{sync_max:thm:families}) to show
sync-maximality for a variety of automata families (among them one for which this was 
previously only conjectured).

Since in the proof of \cref{sync_max:thm:circular_sync_max} we have applied \cref{syncmax:classes:thm:circular} 
when considering certain words as letters, the bounds from \cref{sync_max:prop:don}
do not hold (but they transfer when multiplied by the maximum of the lengths
of these words).

Recall that, by \cref{cor:char-bin-compl-reach}, for binary automata with at least three states,
complete reachability is equivalent to sync-maximality.
However, Complete distinguishability in itself does not imply complete reachability
for binary circular automata. For example, \cref{sync_max:fig:compl_dist_not_compl_reach:1} 
gave a binary circular automaton that is completely distinguishable, but not completely reachable.

\Cref{sync_max:cor:circular} can be used to construct sync-maximal circular automata. We give a few
examples in the following statement. See also \cref{syncmax:fig:binary-circular-sync-maximal} which
depicts some cases from this statement.

\begin{figure}[htbp]
	\centering
	\begin{tikzpicture}
		
		\matrix (m) [matrix of nodes,
		column sep=1.5cm,
		row sep=0.01cm] {
			% nodes={anchor=center}
			
			\begin{tikzpicture}[inner sep=2pt, minimum size=0pt, scale=0.65, transform shape, ->, >=stealth]
				\foreach \i in {0,...,6} {
					\node[state] (q\i) at ({360/7 * \i}:3cm) {$\i$};
				}
				
				\foreach \i [evaluate=\i as \j using {int(mod(\i+1,7))}] in {0,...,6} {
					\path[->] (q\i) edge[bend right=15, dashed] (q\j);
				}
				
				\path[->] (q0) edge (q2);
				\draw (q3) edge (q0); % choose odd i=3				
			\end{tikzpicture}
			
			&
			%	\begin{tikzpicture}[scale=0.75, transform shape, ->, >=stealth, shorten >=1pt, shorten <=1pt]
				\begin{tikzpicture}[scale=0.65, transform shape, ->, >=stealth, node distance=2.2cm, on grid, auto]
					\foreach \i in {0,...,6} {
						\node[state] (q\i) at ({360/7 * \i}:3cm) {$\i$};
					}
						
					\draw (q0) edge[dashed] (q1);
					\draw (q1) edge[bend left=8, dashed] (q2);
					\draw (q2) edge[bend left=8, dashed] (q3);
					\draw (q3) edge[dashed] (q4);
					\draw (q4) edge[bend left=8, dashed] (q5);
					\draw (q5) edge[bend left=8, dashed] (q6);
					\draw (q6) edge[bend left=8, dashed] (q0);
					
					\draw (q0) edge[bend left=20] (q3);
					\draw (q3) edge[bend right=10] (q1); % odd i=3
					
					\foreach \i in {1,2,4,5,6} {
						\pgfmathtruncatemacro{\j}{mod(\i+1,7)}
						\draw (q\i) edge[bend right=8]  (q\j);
					}
				\end{tikzpicture}	
			\\			
			{Automaton as in \crefitem{prop:more_families}{prop:more_families:1}}
			&
			{Automaton as in \crefitem{prop:more_families}{prop:more_families:2}}
			\\
			
			\begin{tikzpicture}[inner sep=5pt, minimum size=0pt, scale=0.65, transform shape, ->, >=stealth, node distance=2.2cm, on grid, auto]
				\foreach \i in {0,...,7} {
					\node[state] (q\i) at ({360/8 * \i}:3cm) {$\i$};
				}
				
				\foreach \i [evaluate=\i as \j using {int(mod(\i+1,8))}] in {0,...,7} {
					\draw (q\i) edge[bend right=15, dashed] (q\j);
				}
				
				\draw (q0) edge (q3);
				\draw (q5) edge (q0);
				
			\end{tikzpicture}
			&
			\begin{tikzpicture}[inner sep=5pt, minimum size=0pt, scale=0.65, transform shape, ->, >=stealth, node distance=2.2cm, on grid, auto]
				\foreach \i in {0,...,7} {
					\node[state] (q\i) at ({360/8 * \i}:3cm) {$\i$};
				}
				
				\foreach \i [evaluate=\i as \j using {int(mod(\i+1,8))}] in {0,...,7} {
					\draw (q\i) edge[bend right=15, dashed] (q\j);
				}
				
				\draw (q0) edge[bend left=10]  (q1);	
				
				\draw (q5) edge[bend left=15]  (q6);
				\draw (q6) edge[bend left=15]  (q7);
				\draw (q7) edge[bend right=30] (q5);
			\end{tikzpicture}
			\\
			
			{Automaton as in \crefitem{prop:more_families}{prop:more_families:3}}
			&
			{Automaton as in \crefitem{prop:more_families}{prop:more_families:4}}
			\\
		};		
	\end{tikzpicture}
	\caption[Example sync-maximal circular automata]{Solid and dashed edges show the actions 
		of~$a$ and, resp., $b$; self-loop omitted.}
	\label{syncmax:fig:binary-circular-sync-maximal}
\end{figure}




\begin{proposition}
	\label{prop:more_families}
	Let $\semiautomat{A}{Q}{\{a,b\}}{\delta}$
	with $Q = [n]$ and $\delta(i, b) = (i + 1) \bmod n$.
	Then~$\mathcal A$ is sync-maximal if~$a$ is defined in the following way:
	\begin{enumerate}
		\item \label{prop:more_families:1} $\delta(0, a) = 2$, precisely one odd $i \in \{1,3,\ldots,n-2\}$
		such that $\delta(i, a) = 0$ and $\delta(j, a) = j$ for $j \in [n] \setminus\{0,i\}$ with $n > 2$ being odd,
		
		\item \label{prop:more_families:2} $\delta(0, a) = 3$, precisely one odd $i \in \{1,3,\ldots,n-2\}$
		such that $\delta(i, a) = 1$ and $\delta(j, a) = (j+1) \bmod n$ for $j \in [n] \setminus\{0,i\}$ with $n > 2$ being odd,
		
		\item \label{prop:more_families:3} $\delta(0, a) = 3$, precisely one $i \in \{1,3,\ldots,n-2\}$
		not divisible by $3$ such that $\delta(i, a) = 0$ and $\delta(j, a) = j$ for $j \in [n] \setminus\{0,i\}$ where $n \not\equiv 0 \pmod{3}$ with $n > 3$,
		
		\item \label{prop:more_families:4} $\delta(0,a) = 1$, $\delta(i, a) = i$ for $i \in \{1,\ldots,\lfloor n / 2 \rfloor\}$,
		and $a$ acts injectively on $\{\lfloor n/2 \rfloor + 1, \ldots, n - 1 \}$
		i.e., $\delta(j, a) = \delta(j', a)$ implies $j = j'$ for $j,j' \in \{\lfloor n/2 \rfloor + 1, \ldots, n - 1 \}$.
	\end{enumerate}
\end{proposition}
\begin{proof}
	All these automata are circular with~$a$ having defect one.
	Then the claims follow from \cref{sync_max:cor:circular} due to the following 
	reasons:
	
	\ref{prop:more_families:1}: We have $\delta(0, a) = \delta(0, b^2a)$ and $n$ is odd.
	
	\ref{prop:more_families:2}: We have $\delta(0, a) = \delta(0, b^2a)$ and $n$ is odd.
	
	\ref{prop:more_families:3}: We have $\delta(0, a) = \delta(0, b^3a)$ and $n$ is not a multiple of~$3$.
	
	\ref{prop:more_families:4}: We have $\delta(0, a) = \delta(0, ba)$.
\end{proof}

As the next example shows, \cref{syncmax:classes:thm:circular} applies to a wider 
class of automata than \cref{sync_max:thm:circular_sync_max}.

\begin{example} 
	\label{ex:thm_not_usable}
	Consider the automaton from \cref{ex:I_vs_Gamma1} (see \cref{fig:thm_not_usable}).
	Let~$w$ be a word of defect one. As can be read off
	from  Table~\ref{tab:rank_3_words},
	for $d > 0$ and $q \in [4]$, 
	if 
	$\applytrans{\delta}{q}{w} = \applytrans{\delta}{\applytrans{\delta}{q}{b^d}}{w}$,
	then $d$ is a multiple of two.
	So \cref{sync_max:thm:circular_sync_max}
	does not apply here.
	We have  $\excl{a} = \{0\}$, $\dupl{a} = \{1\}$,
	and $\delta(0, a) = 1 = \delta(0, b)$.
	Thus \cref{sync_max:prop:don} yields that the automaton
	is completely reachable (which also can be seen in \cref{fig:thm_not_usable}).
\end{example}

Next, we give examples of completely reachable binary automata for which complete reachability does not follow directly from \cref{sync_max:prop:don}. In \cref{syncmax:ex:circular-compl-reach-one}, by considering certain words as letters, it becomes possible to apply \cref{sync_max:prop:don}. However, it is unclear whether such an argument can also be applied in \cref{syncmax:ex:circular-compl-reach-two}.

\begin{example} \label{syncmax:ex:circular-compl-reach-one}
	Let $\semiautomat{A}{Q}{\Sigma}{\delta}$
	where $Q = [n]$ for some $n \ge 3$
	and $\delta(q, b) = (q + 1) \bmod n$ for $q \in Q$,
	$\delta(0, a) = 2$, and $\delta(q, a) = q + 1$
	for $q \in \{1,2,\ldots,n-2\}$
	and $\delta(n-1,a) = 1$.
	For even~$n$, \cref{sync_max:prop:don}
	is not applicable.
	However, we have	
	$\delta(0, a^{n}) = 1$,	
	$a^{n}$ has defect one,
	and $0$ and~$1$ are the excluded and duplicate
	state of~$a^{n}$, respectively. Thus, regarding $a^{n}$
	as an alphabet letter, from \cref{sync_max:prop:don}
	we get that~$\mathcal A$ is completely reachable.
\end{example}

\begin{example} \label{syncmax:ex:circular-compl-reach-two}
	Let $\semiautomat{A}{Q}{\Sigma}{\delta}$
	where $Q = [9]$, $\delta(q, b) = (q + 1) \bmod 9$ for $q \in Q$
	and $\delta(0, a) = 3$, $\delta(a, 6) = 7$, $\delta(7, a) = 6$,
	and $\delta(q, a) = q$ for $q \in Q \setminus\{0,6,7\}$, see \cref{syncmax:fig:example-cr-circular}.
	Next, we show that~$\mathcal A$ is completely reachable.
	Since $\excl{a} = \delta(\dupl{a}, b^3)$,
	this does not follow from \cref{sync_max:prop:don},
	and it is also not clear to us how to find words
	acting in a way to apply \cref{sync_max:prop:don}.
	Next, let $S \subseteq Q$ be a non-empty subset of~$Q$.
	We show, by induction on $|Q \setminus S|$,
	that~$S$ is reachable in $\mathcal A$.
	The claim holds for $S = Q$.
	Next, let $S \ne Q$.	
	
	\begin{enumerate}
	\item Suppose we have $\delta(p, b^3) \in S$ for some $p \in Q \setminus S$.
	
	In this case, there exists $i \ge 0$ such that
	$0 \notin \delta^{-1}(S, b^i)$ and $3 \in \delta^{-1}(S, b^i)$.
	This gives $|\delta^{-1}(S, ab^i)| = |S| + 1$
	and $\delta( \delta^{-1}(S, ab^i), ab^i ) = S$.
	Inductively, we have $\delta(Q, u) = \delta^{-1}(S, ab^i)$
	for some $u \in \Sigma^*$.
	Then $\delta(Q, uab^i) = S$.	
	
	\item Otherwise, each residue class modulo three is either contained in $S$ or in its complement.
		
	Let $q \in Q \setminus S$ be such that $\delta(q, b) \in S$.
	Then we have $6 \notin \delta^{-1}(S, b^i)$ and $7 \in \delta^{-1}(S, b^i)$
	for some $i \ge 0$. 
	Since $6 \notin \delta^{-1}(S, b^i)$
	we also find that $0$ and $3$ are not contained in $\delta^{-1}(S, b^i)$.
	For $T = \delta^{-1}(S, ab^i)$,
	this gives $T = ( \delta^{-1}(S, b^i) \setminus \{ 7  \} ) \cup \{ 6\}$
	and $\delta(T, ab^i) = S$.
	We have $0 \notin \delta^{-1}(T, b^3)$
	since $3 \notin T$ and $3 \in \delta^{-1}(T, b^3)$
	since $6 \in T$.
	We also have $7 \in \delta^{-1}(T, b^3)$ since $1 \in T$
	and $6 \notin \delta^{-1}(T, b^3)$ since $0 \notin T$.
	Thus we find
	\[
	\delta^{-1}(T, ab^3) = ( \delta^{-1}(T, b^3) \setminus \{ 7 \} ) \cup \{ 0,6 \}
	\]
	and $\delta(\delta^{-1}(T, ab^3), ab^3) = T$.
	Since $|\delta^{-1}(T, ab^3)| = |T| + 1 = |S| + 1$, inductively,
	we have $\delta^{-1}(T, ab^3) = \delta(Q, u)$ for some $u \in \Sigma^*$.
	Then we find $\delta(Q, uab^3ab^i) = S$.	
	\end{enumerate}
	
	\begin{figure}[htbp]
		\centering
		\begin{tikzpicture} %[->, >=Stealth, semithick]
			
			\foreach \i in {0,...,8} {
				\node[state] (q\i) at ({40*\i}:3cm) {$\i$};
			}
			
			\foreach \i/\j in {0/1,1/2,2/3,3/4,4/5,5/6,6/7,7/8,8/0} {
				\draw (q\i) edge[dashed] (q\j);
			}
			
			\draw (q0) to (q3);
			\draw (q6) to[bend left=25] (q7);
			\draw (q7) to[bend left=25] (q6);
		\end{tikzpicture}
		\caption[Example completely reachable circular automaton]{Automaton from \cref{syncmax:ex:circular-compl-reach-two}. Solid and dashed edges show the actions 
			of~$a$ and, resp., $b$; self-loop omitted.}
		\label{syncmax:fig:example-cr-circular}
	\end{figure}
\end{example}


\section{Binary Automata}
\label{sync_max:sec:binary}

%!TEX root = ./../../thesis-trier.tex

In \cref{cor:char-bin-compl-reach} 
we have shown that for binary automaton with at least three states,
sync-maximality and complete reachability are equivalent properties.
Here, we show that a binary completely reachable automaton does not admit
any non-trivial congruences. We also show that a binary circular automaton
with a letter of defect one does not admit a non-trivial congruence if and only if
it is synchronizing.

Casas and Volkov~\cite[Thm.~2]{DBLP:conf/latin/CasasV22}
gave the following characterization
of binary completely reachable automata as certain circular automata 
(the authors formulated it differently, assuming a standardized form and identifying
the states with residue classes, that we briefly explain in \cref{syncmax:rem:standardized} below).

\begin{theorem}[{Casas and Volkov~\cite[Thm.~2]{DBLP:conf/latin/CasasV22}}]
	\label{sync_max:thm:casas-volkov-cr-binary}
	An automaton $\semiautomat{A}{Q}{\{a,b\}}{\delta}$ is completely reachable
	if and only if it equals the automaton from \cref{prelim:aut:fig:autnomorph} or
	\begin{enumerate}
		\item one letter, say $a$, has defect one with $\excl{a} = \{q_e\}$,
		
		\item the other letter $b$ cyclically permutes the states, and
		
		\item there does not exist $s \in \delta^{-1}(\dupl{a}, a)$ such that 
		$\applytrans{\delta}{[s]_{\sim}}{a} \subseteq [q_e]_{\sim}$
		for some non-trivial $\langle \fword{\delta}{b} \rangle_G$-congruence $\sim$ on $Q$.
	\end{enumerate}
\end{theorem}

\begin{remark}
	\Cref{sync_max:thm:casas-volkov-cr-binary} easily implies that the automaton
	from \cref{syncmax:ex:circular-compl-reach-two} is completely reachable
	since,  by \cref{lem:congruence_on_cyclic_group},
	its only non-trivial $\langle \fword{\delta}{b} \rangle_G$-congruence~$\sim$
	is the congruence modulo three.
	The excluded state of~$a$ is~$0$ and we have $[0]_{\sim} = \{0,3,6\}$
	and $\delta^{-1}(\dupl{a}, a) = \{0,3\}$.
	Now, $\delta([0]_{\sim}, a) = \{0,7\}$, which is not contained in $[0]_{\sim}$.
	Thus, the automaton admits no non-trivial congruences as in \cref{sync_max:thm:casas-volkov-cr-binary}
	and so it is completely reachable.	 
\end{remark}

\begin{remark}
	We shortly sketch how complete reachability follows
	from the assumptions in \cref{sync_max:prop:don} for binary automata.
	By \cref{lem:congruence_on_cyclic_group}, there exists a divisor $r$ of~$n$
	such that $q_1 \sim q_2$ if and only if 
	$\delta(q_1, b^i) = q_2$ for some multiple $i > 0$ of~$r$.
	Thus, if $\dupl{a}= \applytrans{\delta}{\excl{a}}{b^d}$
	for some $d$ coprime to~$n$, assuming the duplicate and the excluded
	state are congruent implies $r = 1$.
	Hence, the duplicate state and the excluded state of~$a$ are never
	congruent for a non-trivial congruence.
	In \cref{sync_max:thm:casas-volkov-cr-binary}
	we have $\delta(s, a) \sim q_e$ and so the duplicate state and
	the excluded state are congruent.
	Thus, \cref{sync_max:thm:casas-volkov-cr-binary} yields 
	that $\mathcal A$ is completely reachable.
\end{remark}

Originally, Casas and Volkov formulated \cref{sync_max:thm:casas-volkov-cr-binary} 
differently, as we explain next.

\IN{standardized automaton}
\begin{remark} \label{syncmax:rem:standardized}
	Let $\semiautomat{A}{Q}{\{a,b\}}{\delta}$
	be an automaton with $n$ states where $a$ has defect one, $b$ cyclically permutes the states, $s,t \in Q$ are distinct with $\delta(s, a) = \delta(t, a)$, and $\excl{a} = \{q_e\}$.
	Let $0 \le i,j < n$ be such that $\delta(q_e, b^i) = s$
	and $\delta(s, b^j) = q_e$. The words $u = b^ia$ and $v = ab^j$ each
	permute their images and have defect one. Treating them as letters in two
	new automata over the alphabets $\{u,b\}$ and $\{v,b\}$,
	we obtain automata whose transformation monoids equal $\TS{\mathcal A}$;
	so each of them is completely reachable precisely if $\mathcal A$
	is completely reachable. Thus, we can assume $a$ permutes its image.	
	Casas and Volkov~\cite{DBLP:conf/latin/CasasV22} called $\mathcal A$ \emph{standardized}
	if it has this form. (Actually, they also numbered the states from $0$ to~$n - 1$,
	starting with $q_e = 0$, in such a way that applying~$b$ adds one modulo~$n$
	to each state. In this setting, congruences correspond to subgroups and their
	cosets of the cyclic group generated by the permutation induced by~$b$.)
	They showed that~$\mathcal A$ in standardized form is not completely reachable
	if and only if~$a$ preserves the congruence class of $q_e$ for some
	non-trivial congruence.
	Translating this condition to the general (not necessarily standardized) case,
	we obtain \cref{sync_max:thm:casas-volkov-cr-binary} 
	since $\delta([s]_{\sim}, a) \subseteq [q_e]_{\sim}$
	is equivalent to each of the conditions
	that $u$ preserves $[q_e]_{\sim}$ and that~$v$ preserves $[s]_{\sim}$.
\end{remark}

Since $G = \langle \fword{\delta}{b} \rangle_G$ in \cref{sync_max:thm:casas-volkov-cr-binary} 
is transitive, the congruence classes for $\sim$ all have the same size. The next
statement investigates the condition $\applytrans{\delta}{[s]_{\sim}}{a} \subseteq [q_e]_{\sim}$
for general equivalence relation (not necessarily congruences) all of whose
classes have the same size. One of its consequences is that, for the specified condition, it does
not matter which state from $\delta^{-1}(\dupl{a}, a)$ we pick. In fact,
it implies that it can be reformulated to exclude existence of any $q \in Q$
such that $\delta([q]_{\sim}, a) \subseteq [q_e]_{\sim}$ since then
both states from $\delta^{-1}(\dupl{a}, a)$ are mapped into $[q_e]_{\sim}$ by~$a$.

\begin{proposition} \label{syncmax:prop:equivalence-defect-one}
	Let $\semiautomat{A}{Q}{\Sigma}{\delta}$
	with a non-trivial equivalence relation~$\sim$
	all of whose equivalence classes have the same size
	and a letter~$a \in \Sigma$ of defect one
	and $s,t, q_e, q \in Q$ such that $\{s,t\} = \delta^{-1}(\dupl{a}, a)$
	and $\excl{a} = \{q_e\}$.
	Suppose $\delta([q]_{\sim}, a) \subseteq [q_e]_{\sim}$.
	Then we have the following:	
	\begin{enumerate}
		\item \label{syncmax:prop:equivalence-defect-one:1} $\delta([q]_{\sim}, a) = [q_e]_{\sim} \setminus \{q_e\}$.
		
		\item \label{syncmax:prop:equivalence-defect-one:2} $s \sim q$ and $s \sim t$ 
		(and so $\dupl{a} \subseteq [q_e]_{\sim}$).
		
		\item \label{syncmax:prop:equivalence-defect-one:3} For all 
		$p \in Q$ with $p \not\sim s$ we have $\delta(p, a) \not\sim q_e$.
	\end{enumerate}
\end{proposition}
\begin{proof}
	\ref{syncmax:prop:equivalence-defect-one:1}: Since $q_e$ is the excluded state of~$a$,
	we have $\delta([q]_{\sim}, a) \subseteq [q_e]_{\sim} \setminus \{q_e\}$.
	All equivalence classes of~$\sim$ have the same size.
	Assume that $\delta([q]_{\sim}, a)$ is a proper subset of $[q_e]_{\sim} \setminus\{q_e\}$.
	Then we either have three distinct states from $[q]_{\sim}$ mapped by~$a$ to the same state,
	or there are two disjoint pairs of states from $[q]_{\sim}$
	where each pair is mapped to a common state. Both of these cases
	imply that~$a$ has defect at least two, contradicting the assumption
	that~$a$ has defect one. Thus, we must have $\delta([q]_{\sim}, a) = [q_e]_{\sim} \setminus \{q_e\}$.
	
	\ref{syncmax:prop:equivalence-defect-one:2}: Since $\delta([q]_{\sim}, a) = [q_e]_{\sim} \setminus \{q_e\}$
	we have $| \delta([q]_{\sim}, a) | = | [q]_{\sim} | - 1$.
	Thus two states from $[q]_{\sim}$ are mapped to one state by~$a$.
	Since~$a$ has defect one, these two states from $[q]_{\sim}$ must be $s$ and~$t$
	and so $q \sim s$ and $s \sim t$.
	
	\ref{syncmax:prop:equivalence-defect-one:3}: 
	Let $p \in Q$ with $p \not\sim s$. 
	Assume $\delta(p, a) = p'$
	for some $p' \in Q$ with $p' \sim q_e$.
	Since~$p'$ is in the image of~$a$, we have $p' \ne q_e$ and thus $\delta(r, a) = p'$ for some $r \in Q$
	with $r \sim q$ by \cref{syncmax:prop:equivalence-defect-one:1}.
	Thus we also have $r \sim s$ by \cref{syncmax:prop:equivalence-defect-one:2}
	and so $r \ne p$.
	However, since~$a$ has defect one, only $s$ and~$t$ can be mapped
	to the same state. So, this is not possible.	
\end{proof}

Next, we list some properties for equivalence class that are congruences.

\begin{proposition} \label{syncmax:prop:congruences-defect-one}
	Let $\semiautomat{A}{Q}{\Sigma}{\delta}$
	be an automaton with a letter~$a \in \Sigma$ of defect one
	and $s,t \in Q$ such that $\delta^{-1}(\dupl{a}, a) = \{s,t\}$.
	Suppose $\sim$ is a non-trivial $\TS{\mathcal A}$-congruence all of whose
	congruence classes have the same size. Then the following holds:
	\begin{enumerate}
		\item \label{syncmax:prop:congruences-defect-one:1} We have $\delta([p]_{\sim}, a) = [\delta(p, a)]_{\sim}$ for each $p \in Q$ with $p \not\sim s$.
		
		\item \label{syncmax:prop:congruences-defect-one:2} If $\delta(p, a) \sim \delta(q, a)$ for 
		$p,q \in Q$, then $p \sim q$.
		
		%\item \label{syncmax:prop:congruences-defect-one:3} $s \sim t$.		
		
		\item \label{syncmax:prop:congruences-defect-one:4} $\delta([s]_{\sim}, a) \subseteq [q_e]_{\sim}$ for the 
		state $q_e \in Q$ with $\excl{a} = \{q_e\}$.		
	\end{enumerate}
\end{proposition}
\begin{proof}
	\ref{syncmax:prop:congruences-defect-one:1}: Since $\sim$ is a congruence, we have 
	$\delta([p]_{\sim}, a) \subseteq [\delta(p, a)]_{\sim}$.
	The letter~$a$ acts injectively on $Q \setminus \{s\}$.
	We have $[p]_{\sim} \subseteq Q \setminus \{s\}$.
	Thus, since $[p]_{\sim}$ and $[\delta(p, a)]_{\sim}$ have the same size,
	we conclude that $\delta([p]_{\sim}, a) = [\delta(p, a)]_{\sim}$.
	
	\ref{syncmax:prop:congruences-defect-one:2}: Let $\delta(p, a) \sim \delta(q, a)$ for $p,q \in Q$.
	Assume $p \not\sim q$. Without loss of generality, let $p \not\sim s$.
	By \cref{syncmax:prop:congruences-defect-one:1} 
	we have $\delta([p]_{\sim}, a) = [\delta(p, a)]_{\sim}$.
	Consider a state $q_1 \ne q$ with $q_1 \sim q$.
	We have $\delta(q_1, a) \sim \delta(q, a)$.
	Since $\delta([p]_{\sim}, a) = [\delta(p, a)]_{\sim}$,
	there exist $p_1, p_2 \in [p]_{\sim}$ such that $\delta(p_1, a) = \delta(q_1, a)$
	and $\delta(p_1, a) = \delta(q, a)$.
	A letter of defect one maps precisely two states to the same state.
	Thus, since~$p_1, q_1, q$ are pairwise distinct,
	we find that~$a$ has not defect one, contradicting the choice of~$a$.\footnote{Alternatively, 
		$\delta([p]_{\sim}, a) = [\delta(p, a)]_{\sim}$
		and $\delta([q]_{\sim}, a) = [\delta(q, a)]_{\sim}$
		imply $\delta([q]_{\sim}, a) \subseteq \dupl{a}$.
		This yields $\dupl{a} \subseteq [\delta(p, a)]_{\sim}$
		and since all states from $[q]_{\sim}$ are mapped to $\dupl{a}$
		and at least one state from~$[p]_{\sim}$,
		the letter has not defect one.}
	Thus, we must have $p \sim q$.
	
	%\ref{syncmax:prop:congruences-defect-one:3}: Since~$s$ and~$t$ are mapped to the same state, by \cref{syncmax:prop:congruences-defect-one:2} we find $s \sim t$.
	
	\ref{syncmax:prop:congruences-defect-one:4}: 		
	By \cref{syncmax:prop:congruences-defect-one:2}, the letter $a$ acts injectively on the congruence classes.
	Thus, it permutes these classes and for every class there exists some other class
	mapped to it by~$a$.
	By \cref{syncmax:prop:congruences-defect-one:1}, we find $\delta(p, a) \not\sim q_e$ for $p \not\sim s$
	since otherwise not every state in the target congruence class has an
	inverse image in $[p]_{\sim}$.
	So, the only class that can be mapped to $[q_e]_{\sim}$
	is $[s]_{\sim}$.
	
	%Since~$\sim$ is a congruence, we have $\delta([s]_{\sim}, a) \subseteq [\delta(s, a)]_{\sim}$.
	%The sets $[s]_{\sim}$ and $[\delta(s, a)]_{\sim}$ have the same size
	%and~$a$ maps two states from~$[s]_{\sim}$ to the same state by \cref{syncmax:prop:congruences-defect-one:3}, 
	%i.e., it is not injective on this set.
	%This yields that $\delta([s]_{\sim},a)$ is a proper subset of $[\delta(s, a)]_{\sim}$.
	%By \cref{syncmax:prop:congruences-defect-one:2} no state $p \in Q$ with $p \not\sim s$
	%is mapped into $[\delta(s, a)]_{\sim}$.
	%Thus $[\delta(s, a)]_{\sim} \setminus \delta([s]_{\sim},a)$
	%equals the set of excluded states of~$a$,
	%which also gives $\delta(s, a) \sim q_e$.
	%Since $q_e$ is the only excluded state, 
	%we find $\delta([s]_{\sim}, a) = [q_e]_{\sim} \setminus\{q_e\}$.
\end{proof}

\begin{corollary} \label{syncmax:cor:sync-congruences-defect-one}
	Let $\semiautomat{A}{Q}{\{a,b\}}{\delta}$
	be a binary circular automaton with a letter of defect one.
	If $\mathcal A$ is synchronizing,
	then there exists no non-trivial $\TS{\mathcal A}$-congruence.	
\end{corollary}
\begin{proof}
	Let~$a$ be the letter of defect one and~$b$ be the letter cyclically permuting the states.
	Assume there exists a non-trivial $\TS{\mathcal A}$-congruence~$\sim$.
	Then it is also a $G$-congruence for the group generated
	by the permutation induced by~$b$.
	Since~$b$ cyclically permutes the states, and so induces a full cycle on the states,
	$G$ is transitive.
	Thus, all congruence classes have the same size.
	Let $p, q \in Q$ be such that $p \not\sim q$.
	Since~$b$ permutes the states, we have $\delta(p, b) \not\sim \delta(q, b)$.
	By \cref{syncmax:prop:congruences-defect-one}, we also have $\delta(p, a) \not\sim \delta(q, a)$.
	This yields that $\delta(p, u) \not\sim \delta(q, u)$
	for all $u \in \{a,b\}^*$ and so $\mathcal A$ is not synchronizing.	
\end{proof}

In~\cite[Cor.~4]{Volkov24} it was shown that an automaton as in \cref{syncmax:cor:sync-congruences-defect-one}
that does not admit any non-trivial congruences is also synchronizing.

\begin{corollary}
	Let $\semiautomat{A}{Q}{\{a,b\}}{\delta}$
	be a binary completely reachable automaton.
	Then there exists no non-trivial $\TS{\mathcal A}$-congruence.	
\end{corollary}
\begin{proof}
	We can assume $\mathcal A$ has at least three states.
	Suppose $\mathcal A$ is completely reachable.
	By \cref{sync_max:cor2:single-non-permutation} a binary completely reachable automaton
	with more than three states is circular with a letter of defect one.
	Since a completely reachable automaton is synchronizing,
	by \cref{syncmax:cor:sync-congruences-defect-one},  
	$\mathcal A$ admits no non-trivial $\TS{\mathcal A}$-congruences.\footnote{Since
		for a non-trivial congruence~$\sim$
		we would have  $\delta([s]_{\sim}, a) \subseteq [q_e]_{\sim}$ for the 
		state $q_e \in Q$ with $\excl{a} = \{q_e\}$
		and each $s \in \delta^{-1}(\dupl{a}, a)$
		by \cref{syncmax:prop:congruences-defect-one},
		this would also yield that~$\mathcal A$
		does not admit any non-trivial congruences
		if it is completely reachable.}
\end{proof}

Next, in \cref{syncmax:ex:binary-cr-implies-prim}, 
we give an example of a binary circular automaton with a letter of defect
one that does not admit any non-trivial congruence and which is not completely reachable.

\begin{example} \label{syncmax:ex:binary-cr-implies-prim}
	Consider the binary automaton 
	$\semiautomat{A}{Q}{\{a,b\}}{\delta}$ in \cref{sync_max:fig:primitive-not-cr}.
	By \cref{lem:congruence_on_cyclic_group} the only non-trivial $G$-congruence
	for the group~$G$ generated by the full cycle induced by~$b$
	is the congruence modulo three.
	
	\begin{figure}[htb]
		\centering
		\begin{tikzpicture}[shorten >=1pt,->,node distance=4cm]
			
			\foreach \name/\angle/\text in {A/0/0, B/40/1, C/80/2, D/120/3, E/160/4, F/200/5, G/240/6, H/280/7, I/320/8}
			\node[state] (\name) at (\angle:3cm) {$\text$};
			
			\path[->] (A) edge [below] node {$a$} (D);
			\path[->] (B) edge [bend right, above] node {$a,b$} (C);
			\path[->] (C) edge [bend right, below] node {$a$} (B);
			
			\path[->] (A) edge [right] node {$b$} (B)
			(C) edge [above] node {$b$} (D)
			(D) edge [above] node {$b$} (E)
			(E) edge [left] node {$b$} (F)
			(F) edge [left] node {$b$} (G)
			(G) edge [below] node {$b$} (H)
			(H) edge [below] node {$b$} (I)
			(I) edge [right] node {$b$} (A);
		\end{tikzpicture}
		\caption[Not completely reachable and no non-trivial congruences]{An automaton $\mathcal A$ (self-loops omitted) that is not completely reachable and that admits no non-trivial $\TS{\mathcal A}$-congruence, see \cref{syncmax:ex:binary-cr-implies-prim}.}
		\label{sync_max:fig:primitive-not-cr}
	\end{figure}
	
	\begin{enumerate}
		\item $\mathcal A$ admits no non-trivial $\TS{\mathcal A}$-congruence.
		
		We have $1 \equiv 4 \pmod{3}$, $\delta(1, a) = 2$, and $\delta(4, a) = 4$.
		Since $2 \not\equiv 4 \pmod{3}$, congruence classes modulo three do not 
		form a $\TS{\mathcal A}$-congruence.
		Thus there exists no non-trivial $\TS{\mathcal A}$-congruence.	
		
		\item $\mathcal A$ is not completely reachable.
		
		The excluded state of $a$ is $0$, $0$ and~$3$ are mapped
		to the same state, and the letter~$a$ preserves the 
		residue class modulo three of zero. Thus, by \cref{sync_max:thm:casas-volkov-cr-binary}, 
		$\mathcal A$ is completely reachable.
		
	\end{enumerate}
	
	By the above results, it follows that~$\mathcal A$ is synchronizing.
	However, this can also be shown directly in the following way.
	Set $n = 9$. Let $p, q \in Q$ be two distinct states with $\delta(p, b^i) = q$
	for some $0 < i \le 4$.
	Then we have the following cases:
	
	\begin{enumerate}
		\item $i = 2$: Let $u = b^{n - q + 1}$.
		We have $\delta(p, u) = 8$
		and $\delta(q, u) = 1$.
		Then $\delta(p, ua) = 8$
		and $\delta(q, ua) = 2$.
		We have $\delta(8, ba) = \delta(2, ba) = 3$.
		Thus $uaba$ maps both states to~$3$.
		
		\item $i = 1$: Let $u = b^{n - p + 2}$.
		We have $\delta(p, u) = 2$
		and $\delta(q, u) = 3$.
		Thus $\delta(p, ua) = 1$
		and $\delta(q, ua) = 3$.
		We have $\delta(1, b^2) = 3$
		and we have already shown how to map
		two such states to the same state.	
		
		\item $i = 3$: Let $u = b^{n - p}$.
		Then $\delta(p, u) = 0$
		and $\delta(q, u) = 3$.
		Thus $\delta(p, ua) = \delta(q, ua) = 3$.
		
		\item $i = 4$: Let $u = b^{n - p}$.
		We have $\delta(p, u) = 0$
		and $\delta(q, u) = 4$.
		Then $\delta(p, ua) = 3$
		and $\delta(q, ua) = 4$.
		We have $\delta(3, b) = 4$
		and we have already shown how to map 
		two such states to the same state.
	\end{enumerate}
	%
	%	Consider the binary automaton $\semiautomat{A}{Q}{\{a,b\}}{\delta}$ in \cref{sync_max:fig:primitive-not-cr}.
	%It admits no non-trivial $\TS{\mathcal A}$-congruence, which
	%can be easily seen. Also, it is not completely reachable.
	%To show this, arguing by contradiction, assume $\mathcal A$ is completely reachable.
	%Then each of the sets $T_1 = Q \setminus \{0,3,6\}$,
	%$T_2 = Q \setminus \{1,4,7\}$, and $T_3 = Q \setminus \{2,5,8\}$
	%is reachable.
	%For let $u \in \Sigma^*$ be a shortest word
	%reaching one of the sets $T_1, T_2, T_3$.
	%Write $u = vx$ for some $x \in \Sigma$ and $v \in \Sigma^*$.
	%Assume $x = b$. Then $\delta(Q, v)$
	%must also be one of the sets $T_1, T_2, T_3$,
	%contradicting the minimal choice of~$u$.
	%Thus, we have $x = a$
	%and so $\delta(Q, u) \subseteq Q \setminus \{0\}$.
	%Since $0 \in T_2 \cap T_3$,
	%we find $\delta(Q, u) = T_1$.
	%Now, inspecting the action of~$a$, which is the identity on all states distinct from $0,1,2$ 
	%and since $0 \notin T_1$ and $\{1,2\} \subseteq T_1$,
	%this yields $\delta(Q, v) = T_1$,
	%which also contradicts the minimal choice of~$u$.
\end{example}

For a ternary alphabet, complete reachability does not exclude the possibility
of non-trivial $\TS{\mathcal A}$-congruences, see \cref{syncmax:ex:ternary-cr-and-prim}.

\begin{example} \label{syncmax:ex:ternary-cr-and-prim}
	Let $\semiautomat{A}{\{0,1,2,3\}}{\{a,b,c\}}{\delta}$
	where $a$ is a simple idempotent mapping $0$ to~$2$,
	$b$ induces the cycle $(0 ~ 1 ~ 2 ~ 3)$ on the states,
	$\delta(0, c) = \delta(2, c) = 0$, and $\delta(1, c) = \delta(3, c) = 2$.
	
	\begin{enumerate}
		\item The partition $\{0,2\}, \{1,3\}$ of the states induces a $\TS{\mathcal A}$-congruence.
		
		\item $\mathcal A$ is completely reachable.
		
		\item $\{0,1\}$	and $\{2,3\}$ are indistinguishable.
	\end{enumerate} 
	%	In general, in an automaton
	%	having a non-trivial congruence $\sim$ and two non-permutational letters
	%	where one preserves all congruence classes for $\sim$
	%	and the other maps each congruence class to a single state,
	%	no two sets having a state from each congruence class
	%	are distinguishable.
	
	\begin{figure}[htb]
		\centering
		\begin{tikzpicture}[
			->,
			>=stealth',
			shorten >=1pt,
			auto,
			semithick,
			node distance=2.7cm,
			]
			
			\node[state] (0) {$0$};
			\node[state] (1) [right=of 0] {$1$};
			\node[state] (2) [below=of 1] {$2$};
			\node[state] (3) [left=of 2] {$3$};
			
			\path
			(0) edge node[above] {$b$}   (1)
			(1) edge node[right] {$b,c$} (2)
			(2) edge[bend left=10] node[below] {$b$} (3)
			(3) edge node[left]  {$b$} (0);
			
			
			\path (0) edge[bend left=10] node {$a$} (2);
			
			\path
			(2) edge[bend left=10] node[below] {$c$} (0)
			(3) edge[bend left=10] node[above] {$c$} (2);
			
		\end{tikzpicture}
		\caption[Ternary completely reachable automaton with non-trivial congruence]{Automaton 
			(self-loops omitted) that is completely reachable and admits a non-trivial congruence, see \cref{syncmax:ex:ternary-cr-and-prim}.}
		\label{sync_max:fig:ternary-cr-and-prim}
	\end{figure}
\end{example}

%By~\cite[Thm.~2]{DBLP:conf/latin/CasasV22}
%a standardized binary automaton with at least three states is completely reachable
%if and only if %one letter cyclically permutes the states, the other letter has defect one and
%there exists no non-trivial congruence
%such that the congruence class of the excluded state is preserved
%by the letter of defect one. 
%As \cref{syncmax:ex:standardized} shows next, assuming the automaton
%to be standardized is important here.
%
%\begin{example} \label{syncmax:ex:standardized}
%	Consider the automaton $\semiautomat{A}{Q}{\{a,b\}}{\delta}$ in \cref{sync_max:fig:not-prim-not-a-invariant}.
%	The partition $\{0,3,6\}$, $\{1,4,7\}$, $\{2,5,8\}$ induces a $\TS{\mathcal A}$-congruence.
%	Here $\delta(\{0,3,6\}, a) = \{2, 8\}$,
%	$\delta(\{1,4,7\}, a) = \{1,4,7\}$,
%	and $\delta(\{2,5,8\}, a) = \{0, 3, 6\}$.
%	The $\TS{\mathcal A}$-congruence classes containing
%	the excluded state and the two states mapped to the duplicate state
%	are not preserved by~$a$.
%	Thus, it admits no non-trivial $\TS{\mathcal A}$-congruence
%	preserving the congruence class of the excluded state,
%	but it is not completely reachable.
%	This shows that assuming the standardized form for binary automata
%	having a letter cyclically permuting the states and a letter of defect one	
%	in the above characterization derived from~\cite[Thm.~3]{DBLP:conf/latin/CasasV22} is crucial.	
%	
%	\begin{figure}[ht]
	%		\centering
	%		\begin{tikzpicture}[shorten >=1pt,->,node distance=4cm]
		%			
		%			\foreach \name/\angle/\text in {A/0/0, B/40/1, C/80/2, D/120/3, E/160/4, F/200/5, G/240/6, H/280/7, I/320/8}
		%			\node[state] (\name) at (\angle:3cm) {$\text$};
		%			
		%			\path[->] (A) edge [bend left,right] node {$a$} (C);
		%			\path[->] (C) edge [bend right,above] node {$a,b$} (D)
		%			(D) edge [bend right,below] node {$a$} (C)
		%			(F) edge [bend left=18, above] node {$a$} (A)
		%			(G) edge [bend left=40, above] node {$a$} (I)
		%			(I) edge [bend right=80,above] node {$a$} (G);
		%			
		%			\path[->] (A) edge [right] node {$b$} (B)
		%			(B) edge [above] node {$b$} (C)
		%			%(C) edge [above] node {$b$} (D)
		%			(D) edge [above] node {$b$} (E)
		%			(E) edge [left] node {$b$} (F)
		%			(F) edge [left] node {$b$} (G)
		%			(G) edge [below] node {$b$} (H)
		%			(H) edge [below] node {$b$} (I)
		%			(I) edge [right] node {$b$} (A);
		%			
		%			\path[->] (B) edge [loop right] node {$a$} (B)
		%			(E) edge [loop right] node {$a$} (E)
		%			(H) edge [loop above] node {$a$} (H);
		%		\end{tikzpicture}
	%		\caption[Automaton not in standardized form]{Binary automaton $\mathcal A$ not in standardized form
		%			that is not completely reachable and admits no non-trivial $\TS{\mathcal A}$-congruence
		%			such that~$a$ preserves the class of the excluded state, see \cref{syncmax:ex:standardized}.}
	%		\label{sync_max:fig:not-prim-not-a-invariant}
	%	\end{figure}
%\end{example}

\section{Discussion}
\label{sync_max:sec:subclasses:discussion}

%!TEX root = ./../../springer-mono-format/book.tex
%!TEX root = ./../../thesis-trier.tex

\paragraph{Shortest Words Reaching Certain Subsets.} Considering \cref{syncmax:classes:thm:circular}, 
Don~\cite[Conj.~18]{Don16} conjectured
that if a subset of size~$r$ is reachable, then it is reachable by a word of length
at most $n(n - r)$.
Gonze and Jungers~\cite[Prop.~7]{DBLP:journals/jalc/GonzeJ19} refuted this conjecture
by giving a series of $n$-state automata for $n > 6$ having a subset of size
$\lfloor \frac{n}{2} \rfloor$ where the shortest word reaching this subset
has length greater than $\frac{2^n}{n}$. They also noted~\cite[Sect.~2.3]{DBLP:journals/jalc/GonzeJ19}
that truth of this conjecture would have implied the collapse of the polynomial time hierarchy,
which, while still an open problem, is considered to be highly unlikely.
The authors asked whether the conjecture on the length of shortest words
reaching certain subsets might still hold for completely reachable automata.
This was refuted by Zhu~\cite{DBLP:conf/dlt/Zhu24} even for binary automata.
He also proved the bound $n(n-r) + (n-1)$ for binary automata where one letter permutes
its image, and for a subclass of these automata the conjecture was confirmed,
which expanded a result of Casas \& Volkov~\cite{CasasVolkov2024Don} who investigated
this conjecture for binary automata. For circular binary automata having
a letter of defect one that permutes its image the conjecture still remains open~\cite{CasasVolkov2024Don}.
In general, Ferens \& Szyku{\l}a~\cite{FerensSzykula2026} 
proved the bound $2n(n - r) - n \sum_{i=1}^{n - r} \frac{1}{i}$
for completely reachable automata. 

\paragraph{An Erroneous Statement on SI-Automata.} 
In~\cite[Thm.~12]{DBLP:conf/dcfs/Hoffmann22} it was wrongly claimed that a completely reachable SI-automaton is sync-maximal. Figure~\ref{sync_max:fig:SCC_not_sufficient} gives a counterexample.
This wrong conclusion was based on the false statement that
distinguishability of all pairs of distinct $2$-sets
implies complete distinguishability (see the related discussion in~\cref{sync_max:chap_subset_mapping:discussion}).
Here, for automata whose permutational letters generate a transitive group, we proved in \cref{cor:defectatmostone-weaker}
that distinguishability of pairs of distinct $2$-sets, complete distinguishability,
complete reachability, and sync-maximality are equivalent properties.
% implies complete distinguishability,
%and each of both properties is equivalent to complete reachability,
%and complete rea
\Cref{syncmax:cor:SI-2-sets} was also given in~\cite[Lem.~11]{DBLP:conf/dcfs/Hoffmann22}, 
but with  a different (less general) proof.